% =====================================================================================================
% Appendix A: Formal results and proofs.
% v1: module m7_theory (2026-09-23); verification code/analysis/m7_theory/run.py (61 checks).
% v2: module ra5_theory (2026-09-24, referee round 1): generalized wedge (Prop. A8), family token budgets (Prop. A9),
%     model-free sigma* and w (Prop. A10), partial identification of M*(C) (Prop. A11, replaces v1 Prop. A9),
%     global vs local identification on the path (Prop. A1(iii)), quasi-homotheticity (Lemma A4(iv)), DMR sign
%     correction (Prop. A3 and remark), positioning. New and changed claims are verified by
%     code/analysis/ra5_theory/run.py (register: output/tables/ra5_theory_claims.csv).
% Numbering: results are numbered A1, A2, ...; paper/notes/theory_main_text_v2.md maps them to the main text.
% Results A1-A8, Lemmas A1-A5 and Corollaries A1-A5 keep their v1 numbers; A9-A11 are new (v1's A9 is now A11).
% Notation follows paper/notes/model_spec.md; E is irreducible loss and \E (blackboard) is expectation.
% =====================================================================================================

\section{Formal Results and Proofs}\label{app:proofs}

\setcounter{lemma}{0}\setcounter{proposition}{0}\setcounter{corollary}{0}\setcounter{definition}{0}\setcounter{remark}{0}
\renewcommand{\thelemma}{A\arabic{lemma}}
\renewcommand{\theproposition}{A\arabic{proposition}}
\renewcommand{\thecorollary}{A\arabic{corollary}}
\renewcommand{\thedefinition}{A\arabic{definition}}
\renewcommand{\theremark}{A\arabic{remark}}
\renewcommand{\theHlemma}{A\arabic{lemma}}\renewcommand{\theHproposition}{A\arabic{proposition}}
\renewcommand{\theHcorollary}{A\arabic{corollary}}\renewcommand{\theHdefinition}{A\arabic{definition}}\renewcommand{\theHremark}{A\arabic{remark}}
\providecommand{\Var}{\operatorname{Var}}
\providecommand{\Cov}{\operatorname{Cov}}
\providecommand{\E}{\mathbb{E}}

This appendix states and proves the formal results used in the paper. All statements are exact unless an order of
approximation is given. Every algebraic identity below was also verified symbolically and every comparative-static,
bias and information statement numerically (by brute-force optimization, simulation, or direct computation of Fisher
information); the replication package reports each check.

Results here are numbered A1, A2, \dots; the header of each result names the main-text statement it proves. In
brief: Proposition~\ref{prop:duality} (duality) is Lemma~\ref{lem:alloc}; Lemma~\ref{lemma:sigma} ($\sigma<1$ at interior
optima) combines Lemmas~\ref{lem:sigma} and~\ref{lem:soc} with Corollary~\ref{cor:complements};
Proposition~\ref{prop:wedge} (the generalized wedge) is Proposition~\ref{prop:A-wedge}, with Lemmas~\ref{lem:geometry}
and~\ref{lem:ce}, Corollary~\ref{cor:suff} and Proposition~\ref{prop:A-family}; Proposition~\ref{prop:ident}
(identification and information rates) combines Propositions~\ref{prop:fd} and~\ref{prop:A-info} with
Lemma~\ref{lem:geometry}(iv); Proposition~\ref{prop:modelfree} (model-free $\sigma^*$ and $w$) is
Proposition~\ref{prop:A-modelfree}; and Proposition~\ref{prop:pi} (partial identification of $M^*(C)$ and of the wedge) is
Proposition~\ref{prop:A-pi}. Propositions~\ref{prop:A-dmr}--\ref{prop:proxy} and Corollaries~\ref{cor:ceg},
\ref{cor:sahal} and~\ref{cor:hall} support Online Appendix~E (observational production functions and algorithmic
progress). Throughout, $E$ denotes irreducible loss and $\E$ the expectation operator.

\subsection{Setup and notation}\label{app:setup}

\begin{definition}[Technology]\label{def:tech}
A training run uses $N$ parameters and $D$ training tokens (tokens processed, including repetitions) and attains
held-out cross-entropy
\begin{equation}\label{eq:A-tech}
L \;=\; E + e^{-\omega}\Big[A\big(e^{\psi_N}N\big)^{-\alpha} + B\big(e^{\psi_D}D\big)^{-\beta}\Big]e^{\epsilon},
\qquad A,B,\alpha,\beta>0 .
\end{equation}
Lower-case letters are logs: $n=\ln N$, $d=\ln D$, $c=\ln C$ with training compute $C=6ND$, so $c=\ln 6+n+d$; $M\equiv D/N$.
Hicks-neutral productivity $\omega$ and factor-augmenting productivity $(\psi_N,\psi_D)$ are known to the lab when it
chooses inputs; $\epsilon$ is not. Write $u\equiv Ae^{-\alpha(n+\psi_N)}$, $v\equiv Be^{-\beta(d+\psi_D)}$ and
$R\equiv u+v$, so that $L-E=e^{-\omega+\epsilon}R$. (The main text writes $R$ for $L-E$; the two coincide when $\omega=\epsilon=0$.) The log output index is $y\equiv-\ln(L-E)$, and
$F(n,d)\equiv-\ln(Ae^{-\alpha n}+Be^{-\beta d})$. The output elasticities are
$\varepsilon_N\equiv-\partial\ln(L-E)/\partial n$ and $\varepsilon_D\equiv-\partial\ln(L-E)/\partial d$; the wedge is
$w\equiv\varepsilon_N/\varepsilon_D$ and the elasticity share is $s\equiv\varepsilon_N/(\varepsilon_N+\varepsilon_D)$.
Expansion-path objects are $a\equiv\beta/(\alpha+\beta)$, $b\equiv\alpha/(\alpha+\beta)$,
$\gamma\equiv\alpha\beta/(\alpha+\beta)$, $G\equiv(\alpha A/\beta B)^{1/(\alpha+\beta)}$,
$K\equiv\frac{\alpha+\beta}{\beta}AG^{-\alpha}$ and $\sigma^*\equiv2/(2+\alpha+\beta)$. Finally
$\chi\equiv\beta\psi_D-\alpha\psi_N$ (the bias index of productivity) and
$\Omega\equiv\omega+\gamma(\psi_N+\psi_D)$ (productivity in compute-equivalent units).
\end{definition}

\begin{remark}[Exact log form]\label{rem:exact}
Taking logs of \eqref{eq:A-tech} gives, exactly and for any known $E$,
$y=\omega+F(n+\psi_N,d+\psi_D)-\epsilon$, $\varepsilon_N=\alpha u/(u+v)$ and $\varepsilon_D=\beta v/(u+v)$. None of the
three depends on $\omega$, $\epsilon$ or $E$. (Two precision notes on the model specification: the identity is not an
approximation; and ``$\varepsilon_N=\alpha u/R$'' requires $R$ to denote the $\omega$-free sum $u+v$, not the observed
$L-E=e^{-\omega+\epsilon}(u+v)$.)
\end{remark}

\begin{definition}[Generalized family]\label{def:kappa}
The generalized family replaces the bracket in \eqref{eq:A-tech} by
$\big[Ae^{-a_1(n+\psi_N)}+Be^{-b_1(d+\psi_D)}\big]^{\kappa}$ with $a_1,b_1,\kappa>0$. Chinchilla is $\kappa=1$
($a_1=\alpha$, $b_1=\beta$). The joint law of \citet[eq.~1.5]{kaplan2020scaling} is the member with $E=0$,
$a_1=\alpha_N/\alpha_D\approx0.738$, $b_1=1$ and $\kappa=\alpha_D\approx0.103$. For this family write $Q\equiv u+v$
with $u=Ae^{-a_1n}$, $v=Be^{-b_1d}$, $S\equiv a_1+b_1$, $a\equiv b_1/S$, $b\equiv a_1/S$,
$\gamma\equiv\kappa a_1b_1/S$ and $\sigma^*\equiv2/(2+S)$; with $\kappa=1$ these coincide with
Definition~\ref{def:tech}.
\end{definition}

\subsection{The elasticity of substitution and the second-order condition}\label{app:sigma}

\begin{lemma}[Elasticity of substitution; used in Lemma~\ref{lemma:sigma}]\label{lem:sigma}
For the technology \eqref{eq:A-tech} (and for the generalized family with $(a_1,b_1)$ in place of $(\alpha,\beta)$):
\begin{enumerate}
\item[(i)] The Hicks elasticity of substitution between $N$ and $D$ is
\[
\sigma=\frac{\alpha u+\beta v}{\alpha u(1+\beta)+\beta v(1+\alpha)},\qquad\text{equivalently}\qquad \frac1\sigma=1+s\beta+(1-s)\alpha.
\]
\item[(ii)] $\sigma$ depends on $(N,D)$ only through the wedge,
\[
\sigma(w)=\frac{1+w}{1+\alpha+w(1+\beta)},
\]
which is monotone in $w$ between $1/(1+\alpha)$ ($w\to0$) and $1/(1+\beta)$ ($w\to\infty$).
\item[(iii)] On the compute-optimal expansion path ($w=1$), $\sigma=\sigma^*=2/(2+\alpha+\beta)$.
\item[(iv)] $\sigma$ is invariant to any strictly increasing transformation of output; in particular it does not depend
on $E$, $\omega$, $\epsilon$, $\kappa$, or on whether output is measured by loss, bits per byte, or a monotone
benchmark link.
\item[(v)] For the \citet{kaplan2020scaling} law, $\sigma\in[1/2,\,1/1.738]=[0.500,0.575]$ and $\sigma^*=2/3.738=0.535$.
\item[(vi)] With two inputs the Hicks, Allen--Uzawa and Morishima elasticities coincide \citep{blackorby1989real}, so (i)
is all three.
\end{enumerate}
\end{lemma}

\begin{proof}
(i) With output $f(N,D)=-(AN^{-\alpha}+BD^{-\beta})$ in levels, $f_N=\alpha u/N$, $f_D=\beta v/D$,
$f_{NN}=-\alpha(1+\alpha)u/N^2$, $f_{DD}=-\beta(1+\beta)v/D^2$ and $f_{ND}=0$. The Hicks elasticity for two inputs,
$\sigma=-f_Nf_D(Nf_N+Df_D)/\{ND(f_{NN}f_D^2-2f_{ND}f_Nf_D+f_{DD}f_N^2)\}$, then equals
$(\alpha u+\beta v)\alpha u\beta v/\{\alpha u\beta v[\alpha u(1+\beta)+\beta v(1+\alpha)]\}$, which is the stated
expression; dividing numerator and denominator by $\alpha u+\beta v$ gives $1/\sigma$. (ii) Substitute
$\alpha u=w\beta v$. The derivative of $\sigma(w)$ has the sign of $\alpha-\beta$, so $\sigma$ is monotone between the two
limits. (iii) is (ii) at $w=1$. (iv) For $g$ strictly increasing, the first derivatives of $g\circ f$ are $g'$ times
those of $f$, and the second derivatives add terms proportional to $g''f_if_j$; in the denominator these terms sum to
$g''g'^2(f_N^2f_D^2-2f_N^2f_D^2+f_D^2f_N^2)=0$, so numerator and denominator both scale by $g'^3$. Since $L$ is a
strictly increasing transformation of $-R$ (and of $-Q$ in the generalized family) for given $(E,\omega,\epsilon,\kappa)$,
the claim follows. (v) Apply (ii) with $(a_1,b_1)=(0.738,1)$.
\end{proof}

\begin{lemma}[Log coordinates and the second-order condition; Lemma~\ref{lemma:sigma} in the main text]\label{lem:soc}
Let $y=F(n,d)$ be twice continuously differentiable with $F_n,F_d>0$, and define $P\equiv F_nF_d(F_n+F_d)>0$ and
$Q\equiv-(F_{nn}F_d^2-2F_{nd}F_nF_d+F_{dd}F_n^2)$. Then
\begin{enumerate}
\item[(i)] $\sigma=P/(P+Q)$; hence $\sigma\in(0,1)$ if and only if $Q>0$, that is, if and only if the isoquant is
strictly convex in $(\ln N,\ln D)$ coordinates.
\item[(ii)] Consider $\max y$ subject to $n+d=c-\ln6$ (equivalently $\min 6ND$ subject to $y\ge\bar y$). At a point with
$F_n=F_d$, the second derivative of $y$ along the isocost is $F_{nn}-2F_{nd}+F_{dd}=-Q/F_n^2$. The point is a strict
local maximum of output (minimum of compute) if $\sigma\in(0,1)$ and a strict local minimum of output if $Q<0$, which
requires $\sigma>1$ or $\sigma<0$.
\end{enumerate}
\end{lemma}

\begin{proof}
(i) With $f(N,D)=F(\ln N,\ln D)$: $f_N=F_n/N$, $f_D=F_d/D$, $f_{NN}=(F_{nn}-F_n)/N^2$, $f_{DD}=(F_{dd}-F_d)/D^2$ and
$f_{ND}=F_{nd}/(ND)$. Substituting into the Hicks formula, the powers of $N$ and $D$ cancel and
$\sigma=-F_nF_d(F_n+F_d)/[(F_{nn}-F_n)F_d^2-2F_{nd}F_nF_d+(F_{dd}-F_d)F_n^2]=P/(P+Q)$. If $Q>0$ then $P+Q>P>0$ and
$\sigma\in(0,1)$; if $\sigma\in(0,1)$ then $P+Q>P$, so $Q>0$. The curvature of the log-isoquant $d=\phi(n)$ is
$\phi''=Q/F_d^3$, so $Q>0$ is strict convexity. (ii) Along $(n,d)=(n_0-t,d_0+t)$, $d^2y/dt^2=F_{nn}-2F_{nd}+F_{dd}$; at
$F_n=F_d$ the bracket in $Q$ equals $F_n^2(F_{nn}-2F_{nd}+F_{dd})$. If $Q<0$, then $\sigma=P/(P+Q)$ lies outside
$[0,1]$.
\end{proof}

\begin{corollary}[Gross complementarity; Lemma~\ref{lemma:sigma} in the main text]\label{cor:complements}
Under the multiplicative cost $C=6ND$, every interior compute-optimal allocation that satisfies the second-order
condition strictly has $\sigma<1$ at the optimum. If $\sigma>1$ everywhere, compute-optimal training puts all compute on
one input, subject to bounds: for the CES technology $Y=(N^r+D^r)^{1/r}$ with $r\in(0,1)$ ($\sigma=1/(1-r)>1$), output
on the isocost $ND=X$ is $\sqrt X\,(2\cosh(rt))^{1/r}$ with $t=\tfrac12\ln(D/N)$, which is minimized, not maximized,
at the symmetric point. Every member of the generalized family with $a_1,b_1>0$ has $\sigma<1$ everywhere.
\end{corollary}

\begin{proof}
The first two statements are Lemma~\ref{lem:soc}(ii). The CES computation is direct. For the generalized family,
Lemma~\ref{lem:sigma}(ii) gives $\sigma\le1/(1+\min\{a_1,b_1\})<1$.
\end{proof}

\begin{remark}
With a linear cost $p_NN+p_DD$ the second-order condition only requires convex isoquants in levels ($\sigma>0$). With the
multiplicative cost the isocost is a straight line in log coordinates, so the isoquant must be convex in logs, which is
the stronger requirement $\sigma<1$. The implicit level price of parameters, $6D$, rises with the quantity of the other
input; this is why interior optima require complementarity. The empirical content of Corollary~\ref{cor:complements} is
therefore that IsoFLOP profiles are U-shaped in $\ln N$ with an interior minimum; $\sigma^*<1$ is implied by that fact,
not estimated, and every functional form fitted in the paper imposes it. Proposition~\ref{prop:A-modelfree} turns the
curvature of the U into the level of $\sigma^*$.
\end{remark}

\subsection{Duality, allocation and productivity}\label{app:duality}

\begin{lemma}[Cost function and conditional factor demands; Proposition~\ref{prop:duality} in the main text]\label{lem:alloc}
Consider the problem of minimizing $6ND$ subject to $L\le\bar L$, or equivalently of maximizing $y$ at given $c$, for
the technology \eqref{eq:A-tech}.
\begin{enumerate}
\item[(i)] The first-order condition is $\alpha u=\beta v$ ($\varepsilon_N=\varepsilon_D$), and it is sufficient because
$\sigma<1$ everywhere.
\item[(ii)] With $\psi_N=\psi_D=0$: $N^*=G(C/6)^{a}$, $D^*=G^{-1}(C/6)^{b}$ and $R^*(C)=K(C/6)^{-\gamma}$ with
$K=\frac{\alpha+\beta}{\beta}AG^{-\alpha}=(\alpha+\beta)(A/\beta)^{a}(B/\alpha)^{b}$. The cost function is
$C^*(R)=6(K/R)^{1/\gamma}$, with $\partial\ln C^*/\partial\ln R=-1/\gamma$; along the path
$\varepsilon_N=\varepsilon_D=\gamma$.
\item[(iii)] In general,
\begin{align*}
n^*&=\ln G+a(c-\ln6)+a\psi_D-b\psi_N\\
&=\ln G+a(c-\ln 6)+\chi/(\alpha+\beta),\\
d^*&=-\ln G+b(c-\ln6)+b\psi_N-a\psi_D,\\
y^*&=\omega+\gamma(c-\ln6+\psi_N+\psi_D)-\ln K\\
&=\Omega+\gamma(c-\ln6)-\ln K,\\
\ln M^*(C)&=-2\ln G+(b-a)(c-\ln6)-2\chi/(\alpha+\beta).
\end{align*}
\item[(iv)] At given compute, $\omega$ does not enter $(n^*,d^*)$; $\partial n^*/\partial\psi_D=a>0$,
$\partial d^*/\partial\psi_D=-a<0$ and $\partial\ln M^*/\partial\psi_D=-2a<0$: better data raise parameters and lower
both tokens and tokens per parameter. Only the bias index $\chi$ moves the mix; productivity with $\chi=0$
($\alpha\psi_N=\beta\psi_D$) is equivalent to a Hicks-neutral shift $\Delta\omega=\alpha\psi_N$.
\item[(v)] At a given loss target instead of given compute, $c^*-\ln6=(\bar y-\omega+\ln K)/\gamma-\psi_N-\psi_D$, so
$\partial n^*/\partial\psi_D=0$ and $\partial d^*/\partial\psi_D=-1$ (effective data $e^{\psi_D}D$ is held fixed), and
$\partial(d^*-n^*)/\partial\omega=(a-b)/\gamma$, which is nonzero unless $\alpha=\beta$: Hicks-neutral productivity leaves
the expansion path (the mix as a function of compute) unchanged but moves the mix chosen for a given loss target
because the technology is non-homothetic. This co-movement runs entirely through compute: conditional on $c$ the mix
is still flat in $\omega$ (part (iv) and Proposition~\ref{prop:proxy}(ii)); it is the conditioning on compute that rules
out a mix-based proxy for neutral productivity.
\item[(vi)] The sign in (iv) survives inference demand: under the lifetime-compute objective (case (iii)(a) of
Proposition~\ref{prop:A-wedge} with $\lambda=p=1$), $\partial\ln M/\partial\chi=-2/(\alpha+\beta+\theta_T)<0$ at given compute, where
$\theta_T\equiv(T/3D)/(1+T/3D)\in[0,1)$.
\end{enumerate}
\end{lemma}

\begin{proof}
(i) The Lagrangian of $\max\{-\ln(u+v)\}$ subject to $n+d=c-\ln6$ gives $\alpha u/R=\beta v/R$. Sufficiency: the
upper contour sets of $y$ are convex in $(n,d)$ (log-sum-exp is convex), and the constraint is linear.
(ii) Substituting $n=\ln G+a\ln(C/6)$ and $d=-\ln G+b\ln(C/6)$ gives
\[
\frac{\alpha u}{\beta v}=\frac{\alpha A}{\beta B}\,G^{-(\alpha+\beta)}\,(C/6)^{-\alpha a+\beta b}=1,
\]
because $\alpha a=\beta b=\gamma$. At the optimum $v=\alpha u/\beta$, so
$R^*=u(\alpha+\beta)/\beta=\frac{\alpha+\beta}{\beta}AG^{-\alpha}(C/6)^{-\gamma}$. The second expression for $K$ follows by
substituting $G$. Inverting $R^*(C)$ gives $C^*(R)$. On the path $\varepsilon_N=\alpha u/R^*=\alpha\beta/(\alpha+\beta)$.
(iii) Effective inputs $\tilde n=n+\psi_N$ and $\tilde d=d+\psi_D$ satisfy $\tilde n+\tilde d=c-\ln6+\psi_N+\psi_D$ and enter
$R$ exactly as $(n,d)$ do when $\psi=0$. Apply (ii) with effective compute and subtract $\psi_N$ and $\psi_D$. Output is
$y^*=\omega-\ln R^*$ evaluated at effective compute. The form with $\chi$ uses $a\psi_D-b\psi_N=(\beta\psi_D-\alpha\psi_N)/(\alpha+\beta)$.
(iv) Differentiate (iii). If $\alpha\psi_N=\beta\psi_D$ then $u+v=e^{-\alpha\psi_N}(Ae^{-\alpha n}+Be^{-\beta d})$.
(v) Solve $y^*=\bar y$ for $c$ and substitute into (iii).
(vi) By Proposition~\ref{prop:A-wedge} and Corollary~\ref{cor:suff}, the lifetime optimum at compute $c$ satisfies
$\ln M=m_0(c)-2\chi/S+(2/S)\ln(1+\tfrac{T}{3}e^{-d})$ with $S=\alpha+\beta$, $d=\tfrac12(c-\ln6+\ln M)$ and
$m_0(c)=-2\ln G+(b-a)(c-\ln6)$. The implicit function theorem gives
$\partial\ln M/\partial\chi=-(2/S)/(1+\theta_T/S)$.
\end{proof}

\begin{remark}[Chinchilla's approaches]
Lemma~\ref{lem:alloc}(ii) links the three objects that scaling studies estimate: the frontier $R^*(C)$ (a cost function),
the demands $N^*(C),D^*(C)$, and the primal $L(N,D)$. The mapping to \citet{hoffmann2022training} is not one to one.
Their Approach~1 takes the envelope of training curves over FLOPs and reports the argmin $N$ and $D$ as well as the
minimal loss, so it estimates the conditional factor demands as well as the frontier; the frontier fitted through
IsoFLOP minima (their Approach~2 data) is what this paper calls the IsoFLOP-minima frontier. Approach~3 is the primal.
\end{remark}

\subsection{Geometry and the wedge as a sufficient statistic}\label{app:geometry}

\begin{lemma}[Geometry of the technology; used in Propositions~\ref{prop:wedge}, \ref{prop:ident} and~\ref{prop:modelfree}]\label{lem:geometry}
For the generalized family (in effective inputs):
\begin{enumerate}
\item[(i)] $\nabla^2_{(n,d)}\ln R=\kappa\,\dfrac{uv}{Q^2}\begin{pmatrix}a_1^2&-a_1b_1\\-a_1b_1&b_1^2\end{pmatrix}$ has rank
one, and its null vector $(b_1,a_1)$ is the direction of the expansion path $(\partial n^*/\partial c,\partial d^*/\partial c)=(a,b)$.
\item[(ii)] With the transverse coordinate $\tau\equiv a_1n-b_1d$,
$\ln R=\kappa\big[-\tfrac12(a_1n+b_1d)+g(\tau)\big]$ with $g(\tau)=\ln(Ae^{-\tau/2}+Be^{\tau/2})$ and
$g''(\tau)=s_u(1-s_u)\in(0,\tfrac14]$, $s_u=u/Q$. Thus $\ln R$ is affine along every line with direction $(b_1,a_1)$,
and $w$, $s$ and $\sigma$ are constant along such lines.
\item[(iii)] $\ln w=\tau^*-\tau$ with $\tau^*=\ln(a_1A/b_1B)=S\ln G$; the expansion path is the level set
$\{\tau=\tau^*\}$; and for any run with compute $C$ and ratio $M$,
\begin{equation}\label{eq:A-wM}
\ln w=\tfrac{S}{2}\,\ln\!\big(M/M^*(C)\big)=\Big(\frac1{\sigma^*}-1\Big)\ln\!\big(M/M^*(C)\big),
\end{equation}
where $M^*(C)=G^{-2}(C/6)^{b-a}$ is the training-optimal ratio at the run's own compute.
\item[(iv)] (Quasi-homotheticity.) $R(\lambda^{b_1}N,\lambda^{a_1}D)=\lambda^{-\kappa a_1b_1}R(N,D)$ for every $\lambda>0$:
the family is quasi-homogeneous of degree $-\kappa a_1b_1$ with respect to the non-uniform dilation
$(N,D)\mapsto(\lambda^{b_1}N,\lambda^{a_1}D)$, whose orbits are the lines of direction $(b_1,a_1)$ in $(n,d)$. More
generally, let $F$ be $C^2$ with $F_n,F_d>0$ and let $v=(v_n,v_d)$ with $v_n,v_d>0$. Then (a) translation by $tv$ maps the
isoquants of $F$ onto isoquants for every $t$ (the technology is quasi-homothetic with respect to the dilation along $v$)
if and only if (b) $\ln w$ is constant along $v$; and (c) if some strictly increasing transformation $g\circ F$ has $v$ in
the null space of its Hessian everywhere, then (a) and (b) hold. Under (a)--(b) and the multiplicative cost, the set of
compute-optimal points $\{w=1\}$ is a union of orbits; if the optimum at each compute is unique, the expansion path is a
single orbit, a straight line in logs with slope $a=v_n/(v_n+v_d)$, so factor demands are power laws in compute. The
rank-one Hessian in (i) is (c) with $g$ the identity and $v=(b_1,a_1)$: all curvature, and hence all information about
$\sigma$, is transverse to the dilation orbits.
\end{enumerate}
\end{lemma}

\begin{proof}
(i) $\partial\ln Q/\partial n=-a_1u/Q$, so $\partial^2\ln Q/\partial n^2=a_1^2u/Q-a_1^2u^2/Q^2=a_1^2uv/Q^2$; the other
entries follow in the same way, and $\ln R=\kappa\ln Q$. The determinant is zero and
$(a_1^2b_1-a_1b_1a_1,\,-a_1b_1^2+b_1^2a_1)=0$. (ii) Since $-\frac12(a_1n+b_1d)\mp\frac12\tau$ equals $-a_1n$ and $-b_1d$
respectively, $Q=e^{-(a_1n+b_1d)/2}(Ae^{-\tau/2}+Be^{\tau/2})$. Along $(n_0+b_1t,d_0+a_1t)$, $\tau$ is constant and
$a_1n+b_1d$ is affine in $t$. Differentiating $g$ twice gives $pq/(p+q)^2$ with $p=Ae^{-\tau/2}$ and $q=Be^{\tau/2}$, which
is $s_u(1-s_u)$ because $p/(p+q)=u/Q$. (iii) $w=a_1u/(b_1v)=(a_1A/b_1B)e^{-\tau}$. At given compute write
$n=n^*-x/2$ and $d=d^*+x/2$ with $x=\ln(M/M^*(C))$; then $\tau-\tau^*=-(a_1+b_1)x/2$.
(iv) $A(\lambda^{b_1}N)^{-a_1}+B(\lambda^{a_1}D)^{-b_1}=\lambda^{-a_1b_1}(AN^{-a_1}+BD^{-b_1})$; raise to the power $\kappa$.
(a)$\Leftrightarrow$(b): the log-isoquant through a point has slope $dd/dn=-F_n/F_d=-w$. If $w$ is constant along $v$,
the translate of an isoquant solves the same differential equation as the isoquant through the translated point, so the
two coincide; conversely, if translates of isoquants are isoquants, their slopes agree at corresponding points. (c)$\Rightarrow$(b):
if $H_{g\circ F}\,v=0$ everywhere, then $\tfrac{d}{dt}\nabla(g\circ F)(x+tv)=H_{g\circ F}(x+tv)\,v=0$, so the gradient,
and with it the ratio $F_n/F_d$, is constant along $v$. Under (b), $w=1$ at a point implies $w=1$ on its orbit; with a
unique optimum at each compute, the orbit through one optimum contains the optimum at every compute, since each orbit
crosses every isocost $n+d=\text{const}$ exactly once ($v_n+v_d>0$).
\end{proof}

\begin{lemma}[Allocative loss is a function of the wedge; used in Propositions~\ref{prop:wedge} and~\ref{prop:ident}]\label{lem:ce}
Let $C_{\min}(R)$ denote the least training compute $6ND$ that attains reducible loss $R$. For any run with wedge $w$,
\begin{equation}\label{eq:A-ce}
\frac{C}{C_{\min}(R)}=\Big(\frac{a_1+b_1w}{a_1+b_1}\Big)^{S/(a_1b_1)}w^{-1/a_1}
\;\;\overset{\kappa=1}{=}\;\;\Big(\frac{\alpha+\beta w}{\alpha+\beta}\Big)^{1/\gamma}w^{-1/\alpha},
\end{equation}
which depends neither on $(A,B,E,\kappa)$ nor on productivity. If $\alpha=\beta=\rho$,
$C/C_{\min}=\big[\cosh\big(\tfrac\rho2\ln(M/M^*)\big)\big]^{2/\rho}$. To second order,
$\ln(C/C_{\min})=\dfrac{(\ln w)^2}{2S}+O(|\ln w|^3)=\dfrac{\sigma^*}{4(1-\sigma^*)}(\ln w)^2+O(|\ln w|^3)$: the allocative loss is a
Harberger triangle whose area is governed by the elasticity of substitution.
\end{lemma}

\begin{proof}
Let $Q_0=u+v$ at the run. The cheapest bundle producing $Q_0$ has $a_1u^*=b_1v^*$ and $u^*+v^*=Q_0$, so $u^*=aQ_0$ and
$v^*=bQ_0$. Because $u\propto N^{-a_1}$ and $v\propto D^{-b_1}$, $C/C_{\min}=(u^*/u)^{1/a_1}(v^*/v)^{1/b_1}
=(a/s_u)^{1/a_1}\big(b/(1-s_u)\big)^{1/b_1}$. Substituting $s_u=b_1w/(a_1+b_1w)$ gives
$a/s_u=(a_1+b_1w)/(Sw)$ and $b/(1-s_u)=(a_1+b_1w)/S$, and $1/a_1+1/b_1=S/(a_1b_1)$. With $a_1=b_1=\rho$,
\eqref{eq:A-ce} equals $\big[(w^{1/2}+w^{-1/2})/2\big]^{2/\rho}=\cosh(\tfrac12\ln w)^{2/\rho}$ and $\ln w=\rho\ln(M/M^*)$ by
\eqref{eq:A-wM}. For the expansion, let $x=\ln w$: $\ln\big((a_1+b_1e^x)/S\big)=ax+\tfrac12abx^2+O(x^3)$; multiplying by
$S/(a_1b_1)=1/(abS)$ and subtracting $x/a_1=x/(bS)$ leaves $x^2/(2S)$. Finally $S=2(1-\sigma^*)/\sigma^*$.
\end{proof}

\subsection{What on-path data identify}\label{app:onpath}

\begin{definition}[On-path sample]\label{def:onpath}
A sample of runs is on-path if (a) all runs share the technology (generalized family, common parameters); (b)
productivity is Hicks-neutral ($\psi\equiv0$); (c) every run minimizes training compute for its loss ($T=0$, no binding
constraints); and (d) compute varies over an interval (or at least five distinct values).
\end{definition}

\begin{proposition}[What on-path data identify; Proposition~\ref{prop:ident}(i) in the main text]\label{prop:fd}
In an on-path sample:
\begin{enumerate}
\item[(i)] $n_i=\ln G+a(c_i-\ln6)$ and $d_i=-\ln G+b(c_i-\ln6)$. The inputs are deterministic affine functions of compute,
and the choice data reveal $a$ and $G$ (the observed compute-optimal ratio $M^*(C)$).
\item[(ii)] $\E[y\mid c]=\gamma(c-\ln6)-\ln K+\E[\omega\mid c]$. Within this model (without instruments for compute or
panel restrictions on $\omega$), the frontier elasticity $\gamma$ is identified if and only if $\E[\omega\mid c]$ is known up
to an additive constant (for example, $\omega$ mean-independent of $c$, as in a single-lab sweep).
\item[(iii)] Suppose in addition $\E[\omega\mid c]$ is constant and $\kappa=1$. Write $c'\equiv c-\ln6$, $r_N\equiv\alpha a$
and $r_D\equiv\beta b$ for the rates at which the two terms decay along the observed path, and $P\equiv Ae^{-\alpha n_0}$,
$Q\equiv Be^{-\beta d_0}$ (with $K$ normalized by $\E[e^{-\omega+\epsilon}]=1$), so that on the path
$\E[L\mid c]=E+Pe^{-r_Nc'}+Qe^{-r_Dc'}$: a two-rate exponential mixture whose rates coincide at the truth
($r_N=r_D=\gamma$).
\begin{enumerate}
\item[(a)] (Global identification, by functional form.) $E$, $\gamma$ and $K$ are identified, and so are $\alpha=\gamma/a$
and $\beta=\gamma/b$: every parameter with $A,B>0$ that generates the same on-path outcomes has $r_N=r_D=\gamma$. The ratio
$A/B$, and hence the technology's own cost-minimizing ratio, is identified only by imposing that the observed path is
optimal ($G^{\alpha+\beta}=\alpha A/\beta B$); outcomes identify only $P+Q$.
\item[(b)] (No local identification.) The Jacobian of the on-path mean of $\ln L$ with respect to $(E,\ln A,\ln B,\alpha,\beta)$
has rank three; its null space contains the direction that changes $A/B$ at fixed $K$ and a direction with
$\beta^2\delta\alpha+\alpha^2\delta\beta=0$. Because $A/B$ is not identified, more is true: every perturbation that moves
$r_N$ and $r_D$ in opposite directions leaves the on-path mean unchanged to first order (the change is absorbed by
re-splitting $P+Q$), so along each such direction the least-squares or likelihood criterion is of fourth order in the
step. In particular $\sigma^*=2/(2+\alpha+\beta)$ is identified only at second order \emph{whether or not} $\alpha=\beta$;
the gradient of $\sigma^*$ lies in the row space of the Jacobian if and only if $\alpha=\beta$, but that Jacobian is
evaluated at the true $A/B$, which on-path data do not reveal. (If $A/B$ is known, or pinned by imposing the intercept
restriction $G^{\alpha+\beta}=\alpha A/\beta B$, the only first-order-flat direction is the one that keeps the split
$P/Q$ at its true value, and $\sigma^*$ is then first-order identified if and only if $\alpha=\beta$, which is the first
draft's statement; it does not survive estimating $A/B$ from outcomes.) Moreover, the \emph{twin} map
$(\alpha,\beta,A,B)\mapsto(\beta b/a,\ \alpha a/b,\ Be^{-\beta d_0+\beta bn_0/a},\ Ae^{-\alpha n_0+\alpha ad_0/b})$ leaves the
on-path outcome distribution exactly unchanged ($a,b$ are the observed path slopes), so every parameter value off the
ray $\alpha a=\beta b$ has an observationally equivalent twin, whose $\sigma^*$ differs unless $a=b$ (the truth is its own
twin). On the boundary of the parameter
space identification fails altogether: with $A\to0$ ($B\to0$) the on-path distribution converges to the true one for
any $\alpha$ ($\beta$).
\item[(c)] (Inference.) The on-path model is an over-fitted mixture: the spread $|r_N-r_D|$ enters the mean only through a
second-order term, so, by the analogue of the optimal rate for over-fitted mixtures \citep{chen1995optimal}, no estimator of
$(\alpha,\beta)$ or $\sigma^*$ should be expected to converge faster than $n^{-1/4}$ uniformly near the truth. This is a
heuristic argument, not a theorem for this model (a nonlinear regression, not a mixture density); see
\citet{sargan1983identification} on identification without first-order identification and \citet{rotnitzky2000likelihood}
on likelihood inference with a singular information matrix, whose results do not apply directly because $P/Q$ is not
identified at the truth. In a compact parameter space bounded away from $A=0$ and $B=0$, least squares attains this rate
(by simulation, below); its
distribution is non-normal (a mass of estimates with $\hat r_N=\hat r_D$, which converge at $n^{-1/2}$, and a mass with split
rates); Wald and delta-method intervals based on the (near-)singular Hessian are invalid; and the likelihood-ratio
statistic does not have the standard $\chi^2$ limit. Imposing the cross-equation restriction $\beta/(\alpha+\beta)=a$ (the fitted
technology's allocation exponent equals the observed path slope) forces $r_N=r_D$ and restores $n^{-1/2}$ rates. The
precision of on-path estimates of $\sigma^*$ under $\kappa=1$ is therefore the precision of that restriction: functional
form plus optimality.
\end{enumerate}
\item[(iv)] Without the restriction $\kappa=1$: for every $S>0$, the member
\begin{equation}\label{eq:A-family}
\begin{gathered}
a_1=(1-a)S,\qquad b_1=aS,\qquad \kappa_S=\frac{\gamma}{a(1-a)S},\\
A_S=\frac{b_1}{S}K^{1/\kappa_S}G^{a_1},\qquad B_S=\frac{a_1A_S}{b_1G^{S}}
\end{gathered}
\end{equation}
generates exactly the same joint distribution of $(n,d,y)$ as the true technology. Hence $\sigma^*=2/(2+S)$ is not
identified: every value in $(0,1)$ is consistent with the data, whereas by Lemma~\ref{lem:soc} no value $\sigma^*\ge1$ is,
because such technologies make the observed allocations cost maxima.
\end{enumerate}
\end{proposition}

\begin{proof}
(i) is Lemma~\ref{lem:alloc}(ii). (ii) On the path $y=\omega+\gamma(c-\ln6)-\ln K-\epsilon$ with $\E[\epsilon\mid c]=0$;
if $\E[\omega\mid c]=h(c)$ with $h$ unknown, then $\gamma(c-\ln6)+h(c)$ and $\tilde\gamma(c-\ln6)+h(c)+(\gamma-\tilde\gamma)(c-\ln6)$
are indistinguishable. (iii)(a) With $\omega$ independent of $c$, $\E[L\mid c]=E+\tilde Ke^{-\gamma c'}$ with
$\tilde K=K\,\E[e^{-\omega+\epsilon}]$; a three-parameter exponential is identified from an interval of $c$. Suppose
$(\tilde\alpha,\tilde\beta,\tilde A,\tilde B)$ with $\tilde A,\tilde B>0$ fits the same curve on the observed path
$n=n_0+ac'$, $d=d_0+bc'$. Then $\tilde Pe^{-\tilde\alpha ac'}+\tilde Qe^{-\tilde\beta bc'}=\tilde Ke^{-\gamma c'}$ on an
interval with $\tilde P,\tilde Q>0$. Exponentials with distinct rates are linearly independent, so
$\tilde\alpha a=\tilde\beta b=\gamma$, which gives $\tilde\alpha=\gamma/a$ and $\tilde\beta=\gamma/b$. The remaining
condition, $\tilde P+\tilde Q=\tilde K$, is one equation in two unknowns; imposing optimality adds
$G^{\alpha+\beta}=\alpha A/\beta B$.
(b) For the Jacobian, on the path $u=u_0e^{-\gamma c'}$ and $v=v_0e^{-\gamma c'}$ with $u_0/v_0=\beta/\alpha$. The scores
of $\ln L$ are $1/L$ times: $1$ for $E$; $u$ and $v$ for $\ln A$ and $\ln B$; $-(n_0+ac')u$ for $\alpha$; and
$-(d_0+bc')v$ for $\beta$. They span $\{1,e^{-\gamma c'},c'e^{-\gamma c'}\}$, so the rank is three. A null vector must
satisfy $\delta E=0$, $au_0\delta\alpha+bv_0\delta\beta=0$ (the coefficient on $c'e^{-\gamma c'}$), and
$u_0\delta\ln A+v_0\delta\ln B-u_0n_0\delta\alpha-v_0d_0\delta\beta=0$; with $u_0/v_0=\beta/\alpha$ the second condition
is $\beta^2\delta\alpha+\alpha^2\delta\beta=0$, along which $\delta(\alpha+\beta)=\delta\alpha(1-\beta^2/\alpha^2)$. This
local analysis holds $P/Q$ at its true value. Every $(P,Q)$ with $P+Q=P_0+Q_0$ and equal rates generates the true
distribution, and at such a point the first-order change of the mean from $(\delta r_N,\delta r_D)$ is
$-c'e^{-\gamma c'}(P\delta r_N+Q\delta r_D)$, which vanishes for $P/Q=-\delta r_D/\delta r_N>0$ whenever
$\delta r_N\delta r_D<0$. The second-order term, $\tfrac12c'^2e^{-\gamma c'}(P\delta r_N^2+Q\delta r_D^2)$, is not in the span of
$\{1,e^{-\gamma c'},c'e^{-\gamma c'}\}$, so the criterion is exactly of fourth order along such directions. Since
$\delta(\alpha+\beta)=\delta r_N/a+\delta r_D/b$, generic opposite-sign directions move $\sigma^*$ at first order, including
when $a=b$. The twin claim follows by substitution: the twin's two terms are $Qe^{-r_Dc'}$ and $Pe^{-r_Nc'}$ with the roles of
the inputs exchanged. On the boundary, $A=0$ removes the $N$-term whatever $\alpha$ is, and the remaining term fits the true
curve with $\beta=\gamma/b$.
(c) Writing the on-path mean as $E+(P+Q)\int e^{-rc'}\,dG(r)$ with $G$ the two-point distribution on $\{r_N,r_D\}$ with
weights proportional to $(P,Q)$, the data identify the moments of $G$ around $\gamma$ through the coefficients on
$c'^ke^{-\gamma c'}$: the mean from the first-order term and the variance from the second-order term, each at best at rate
$n^{-1/2}$, so the spread $|r_N-r_D|$, and with it $(\alpha,\beta)$ and $\sigma^*$, at best at rate $n^{-1/4}$, as in
\citeauthor{chen1995optimal}'s bound for over-fitted mixtures (a heuristic argument; the simulation below confirms the rate
for least squares). The restriction $\beta/(\alpha+\beta)=a$ is $r_N=r_D$, which leaves the regular model $E+(P+Q)e^{-rc'}$. A Monte Carlo with 40 on-path compute levels confirms the rates: the 75th and
90th percentiles of $|\hat\sigma^*-\sigma^*|$ scale as $n^{-0.24}$ and $n^{-0.27}$ to $n^{-0.28}$ both for
$\alpha\neq\beta$ and for $\alpha=\beta$, and as $n^{-0.50}$ with the restriction imposed (replication package).
(iv) Apply Lemma~\ref{lem:alloc}(ii) to the inner aggregator of member $S$. Its path slope is $b_1/S=a$. Its path
intercept is $(a_1A_S/b_1B_S)^{1/S}=G$ by construction of $B_S$. Its inner frontier is $K_1(C/6)^{-a_1b_1/S}$ with
$K_1=(S/b_1)A_SG^{-a_1}=K^{1/\kappa_S}$, by construction of $A_S$. Hence $R^*=K(C/6)^{-\kappa_Sa_1b_1/S}=K(C/6)^{-\gamma}$.
Inputs and outputs therefore coincide run by run for every $(\omega,\epsilon)$. Values $\sigma^*\ge1$ correspond to
$S\le0$ (with $\kappa<0$) and violate Lemma~\ref{lem:soc}(ii).
\end{proof}

\begin{remark}[Identification in practice]
Proposition~\ref{prop:fd}(iii) concerns on-path designs: compute-optimal ladders and fixed tokens-per-parameter designs.
The Chinchilla data are not such a design; they have transverse IsoFLOP variation, and the large standard errors that
\citet{besiroglu2024chinchilla} report for $A$ and $B$ (125 and 1,293) mostly reflect the parameterization at $N=D=1$
(normalizing at the geometric means of the inputs, as in \citealp{klump2007factor}, removes most of the ill-conditioning).
The distinction that matters is between extrapolation in compute and on-path non-identification. IsoFLOP-minima
estimators (Approach~2, and the frontier through the minima) estimate $(a,G)$ from the argmins of designed profiles, that
is, from outcomes of assigned inputs; with behavioral on-path data the same objects come from choices. The primal
(Approach~3) estimates $(E,A,B,\alpha,\beta)$ from outcomes; the two are linked by $G^{\alpha+\beta}=\alpha A/\beta B$, a
restriction that Hoffmann et al.'s Approach~3 violates \citep{hoffmann2022training}.
\end{remark}

\begin{proposition}[Information rates; Proposition~\ref{prop:ident}(ii) in the main text]\label{prop:A-info}
Consider a sweep with $y_i=F(n_i,d_i;\phi)-\epsilon_i$, $\epsilon_i$ i.i.d.\ $N(0,V_\epsilon)$, $E$ known, and the generalized
family parameterized by $\phi=(S,a,\gamma,\ln G,\ln K)$. For run $i$ let $w_i$ be its wedge,
$\Delta_i\equiv bn_i-ad_i-\ln G=-(\ln w_i)/S$ its transverse deviation (Lemma~\ref{lem:geometry}), and
$\Phi_i\equiv\ln(C_i/C_{\min}(R_i))$ its allocative loss (Lemma~\ref{lem:ce}).
\begin{enumerate}
\item[(i)] (Exact decomposition.) $y_i=\gamma(c_i-\ln6)-\ln K-\gamma\,\Phi_S(\Delta_i)-\epsilon_i$, where
$\Phi_S(\Delta)=\big[\ln(b+ae^{-S\Delta})+aS\Delta\big]/(abS)=\Phi_i\ge0$ and $\Phi_S(0)=\Phi_S'(0)=0$. The curvature
parameter $S$, and hence $\sigma^*$, enters output only through the runs' allocative losses.
\item[(ii)] (Scores.) The scores for the path level and for curvature are
\begin{align*}
\frac{\partial y_i}{\partial\ln G}&=\frac{\gamma(1-w_i)}{b+aw_i}=-\gamma\ln w_i+O\big((\ln w_i)^2\big),\\
\frac{\partial y_i}{\partial S}&=-\gamma\,\frac{\partial\Phi_S(\Delta_i)}{\partial S}=-\frac{\gamma}{S}\,\Phi_i\,[1+O(|\ln w_i|)]
=-\frac{\gamma(\ln w_i)^2}{2S^2}+O(|\ln w_i|^3).
\end{align*}
\item[(iii)] (Efficient information.) Let $\zeta_k$ be the vector of scores for parameter $k$ and $M_{-k}$ the
residual-maker for the other four score vectors. Then
\begin{align*}
\mathcal I^*(\ln G)&=V_\epsilon^{-1}\|M_{-G}\zeta_G\|^2\\
&=(\gamma^2/V_\epsilon)\,\|M_{-G}\ln w\|^2\,[1+O(\textstyle\max_i|\ln w_i|)],\\
\mathcal I^*(S)&=V_\epsilon^{-1}\|M_{-S}\zeta_S\|^2\\
&\le\big(\gamma^2/(S^2V_\epsilon)\big)\textstyle\sum_i\Phi_i^2\,[1+O(\max_i|\ln w_i|)]\\
&=O\big(\textstyle\sum_i(\ln w_i)^4\big).
\end{align*}
If every run lies on the expansion path ($w_i\equiv1$), both informations are zero. If the transverse deviations of a
fixed design are scaled by $v\to0$, then $\mathcal I^*(\ln G)=O(v^2)$ and $\mathcal I^*(\sigma^*)=O(v^4)$, so the
standard error of $\hat\sigma^*$ grows like $v^{-2}$ and that of $\ln\hat M^*$ like $v^{-1}$. Both rates are attained
by factorial designs.
\item[(iv)] (Estimated $E$, noise on $\ln L$.) If $E$ is estimated and the noise is Gaussian on $\ln L$ instead of $y$, the
score for $S$ is multiplied by $R_i/L_i\in(0,1)$, so the bound $\mathcal I^*(S)=O\big(\sum_i(\ln w_i)^4\big)$ still holds;
in the nine-budget, five-offset factorial design the orders $v^4$ (for $S$) and $v^2$ (for $\ln M^*$ at the design's
central compute) are attained (numerical log-log slopes 4.00 and 2.00), and estimating $E$ costs no information about
$S$ to two decimals.
\end{enumerate}
\end{proposition}

\begin{proof}
(i) Because the frontier is a power law, $R=K(C_{\min}(R)/6)^{-\gamma}$ for every $R$, so
$\ln R=\ln K-\gamma(c-\ln6)+\gamma\ln(C/C_{\min}(R))$. By Lemma~\ref{lem:ce} for the generalized family and
Lemma~\ref{lem:geometry}(iii), $\ln(C/C_{\min})=\ln\big((a_1+b_1w)/S\big)/(abS)-\ln w/a_1$ with $w=e^{-S\Delta}$,
$a_1=bS$ and $b_1=aS$; this is $\Phi_S(\Delta)$. It is nonnegative because $C\ge C_{\min}$, and $\Phi_S(0)=0$. Its
derivative is $\Phi_S'(\Delta)=(1-e^{-S\Delta})/(b+ae^{-S\Delta})$, which vanishes at zero.
(ii) Holding $(n,d)$ fixed, $\Delta$ depends on $\ln G$ (with derivative $-1$) but not on $S$. Hence
$\partial y/\partial\ln G=\gamma\Phi_S'(\Delta)=\gamma(1-w)/(b+aw)$. Write $\Phi_S=h(S\Delta)/(abS)$ with
$h(x)=\ln(b+ae^{-x})+ax$; then $\partial\Phi_S/\partial S=[xh'(x)-h(x)]/(abS^2)$, and
$h(x)=\tfrac12abx^2+O(x^3)$, $xh'(x)-h(x)=\tfrac12abx^2+O(x^3)$. So
$\partial\Phi_S/\partial S=\tfrac12\Delta^2+O(|\Delta|^3)=\Phi_S/S+O(|\Delta|^3)$.
(iii) Under Gaussian errors the information matrix is $V_\epsilon^{-1}Z'Z$ with $Z$ the score matrix; the efficient
information for one parameter is the Schur complement $V_\epsilon^{-1}\zeta_k'M_{-k}\zeta_k$. The inequality uses
$\|M_{-S}\zeta_S\|\le\|\zeta_S\|$ and (ii). The orders follow from (ii): $\zeta_G=O(v)$ and $\zeta_S=O(v^2)$.
Attainment: for a design that crosses several compute levels with symmetric transverse offsets, $\Delta_i^2$ is not in
the span of the leading-order nuisance scores $\{1,\,c_i-\ln6-\Phi_i,\,\Delta_i,\,\Delta_i(c_i-\ln6)\}$ and $\Delta_i$ is not in
the span of the others. The numerical exponents are 4.00 and 2.00 for a nine-budget, five-offset factorial design.
(iv) $\ln L=\ln(E+R)$, so $\partial\ln L/\partial S=(R/L)\,\partial\ln R/\partial S=-(R/L)\,\partial y/\partial S$ and the bound
in (iii) applies with the factor $(R_i/L_i)^2\le1$; the attainment claims are numerical (finite-difference scores with $E$
as a sixth parameter).
\end{proof}

\begin{remark}[Relation to the literature, and where the results bind]\label{rem:A-positioning}
Propositions~\ref{prop:fd} and~\ref{prop:A-info} are instances of the classical point that optimizing firms facing
common relative prices generate inputs that lie on a lower-dimensional set \citep{marschak1944random}.
\citet{bond2005adjustment} show that in a Cobb--Douglas model with common input prices and no adjustment costs the input
elasticities are not identified, and that identification requires variation in relative prices, adjustment frictions or
optimization errors; \citet{gandhi2020identification} show that flexible-input elasticities are not identified from the
proxy structure alone. Scaling laws are the cleanest instance: because the FLOP accounting $C\approx6ND$ is common, the
log-cost weights are identical across labs, and without wedges every allocation lies on one ray (Proposition~\ref{prop:fd}(i)).
This is distinct from the functional dependence of \citet{ackerberg2015identification}, which concerns the first stage of
proxy estimators (Proposition~\ref{prop:proxy}(iii)). The escape routes map one to one: relative-price variation
corresponds to heterogeneous wedges (serving demand, data costs, size-dependent training prices; Proposition~\ref{prop:A-wedge});
adjustment frictions to binding memory, latency or data constraints; optimization errors to belief errors, such as
allocations planned with the \citet{kaplan2020scaling} law. Designed IsoFLOP sweeps are a different escape: the
experimenter assigns $(N,D)$ directly, so there is neither a price nor a behavioral response.

The rates in Proposition~\ref{prop:A-info} are the design-theoretic property of estimating a location parameter versus a
curvature (second-order) parameter in nonlinear regression: information about a curvature parameter scales with the
fourth moment of the design's spread \citep{box1959design,kiefer1959optimum}. What is specific here is that the spread is
generated by economic wedges, and that the curvature enters output only through the runs' Farrell allocative losses
(Proposition~\ref{prop:A-info}(i)). Proposition~\ref{prop:fd}(iii) separates identification by functional form, which holds
globally, from local (first-order) identification, which fails, with the non-standard inference that follows. It also
reconciles two statements that look contradictory. \citet{kricheli2026tokens} show that fixed tokens-per-parameter designs
become ill-conditioned as $|\alpha-\beta|\to0$; the first draft of this paper said that on the compute-optimal path
$\sigma^*$ is first-order identified if and only if $\alpha=\beta$. The mechanism is one: the degeneracy arises whenever the two
terms decay at the same rate along the design's ray, $\alpha a'=\beta b'$ for a ray of slope $a'=1-b'$. On a fixed-ratio
ray ($a'=b'=\tfrac12$) that happens if and only if $\alpha=\beta$; on the compute-optimal path it happens always, because
equal decay rates are what compute-optimality means ($\varepsilon_N=\varepsilon_D$ along the path), and then, with $A/B$
estimated, $\sigma^*$ is not first-order identified whether or not $\alpha=\beta$ (the first draft's statement holds only
with $A/B$ known).

Where the results bind: on compute-optimal ladders and fixed tokens-per-parameter designs run by a single lab, and on
releases of the compute-optimal era. Released models since 2023 lie far off the training-optimal path. Their heterogeneous
wedges supply transverse variation, but that variation is generated by the objectives of Proposition~\ref{prop:A-wedge}
and is confounded with factor-biased technology (Proposition~\ref{prop:proxy}(iv)), so it identifies the technology only
jointly with a model of the wedges. (For $\sigma^*$ itself, multiply the information bound for $S$ by
$(dS/d\sigma^*)^2=4/\sigma^{*4}$; the order is unchanged.)
\end{remark}

\subsection{Technical change over time}\label{app:dmr}

\begin{proposition}[What time series of compute-optimal models identify; used in Online Appendix~E]\label{prop:A-dmr}
Suppose compute-optimal models are observed at several dates.
\begin{enumerate}
\item[(i)] (Nonparametric sign result.) Let $y=F(n,d,t)$ be smooth with $F_n,F_d>0$, and let the compute-optimal
allocation be interior at each date. Define the Hicks bias $\mathcal B_t\equiv\partial\ln(F_n/F_d)/\partial t$ at fixed $(n,d)$
(the change in the marginal rate of technical substitution at fixed inputs; $\mathcal B_t>0$ is parameter-using). Then
\begin{equation}\label{eq:A-drift}
\frac{\partial\ln M^*(C,t)}{\partial t}=-\frac{\sigma^*_t}{1-\sigma^*_t}\,\mathcal B_t,\qquad\sigma^*_t\in(0,1).
\end{equation}
Hence the sign of the Hicks bias is identified from compute-optimal models alone: $\mathcal B_t>0$ if and only if the
compute-optimal tokens-per-parameter ratio falls at given compute, and neutrality ($\mathcal B_t=0$) is testable. The magnitude of
$\mathcal B_t$ is not identified without $\sigma^*_t$.
\item[(ii)] (Generalized family with factor augmentation $\psi_{N,t}=g_Nt$, $\psi_{D,t}=g_Dt$.) The frontier shifts as
$\ln K_t=\ln K_0-\gamma(g_N+g_D)t$; the path intercept moves as $\ln G_t=\ln G_0+(ag_D-bg_N)t$; the mix drifts as
$\partial\ln M^*/\partial t=2(bg_N-ag_D)$; and $\mathcal B=b_1g_D-a_1g_N=S(ag_D-bg_N)$, which equals $(1-1/\sigma)(g_N-g_D)$ only
when $a_1=b_1$. The drift of $M^*$ therefore signs $b_1g_D-a_1g_N$ ($\beta g_D-\alpha g_N$ when $\kappa=1$), which is the
sign of $g_D-g_N$ only when $a_1=b_1$. Because $a$ and $\gamma$ are identified at each date (Proposition~\ref{prop:fd}),
$g_N$ and $g_D$ are identified separately, by the family's functional form. $S$, and therefore $\sigma^*$ and $|\mathcal B|$, are not: the family \eqref{eq:A-family} with
$(A_{S,t},B_{S,t})=(A_Se^{-a_1g_Nt},B_Se^{-b_1g_Dt})$ is observationally equivalent at every date for every $S>0$.
\item[(iii)] Hicks neutrality $\Leftrightarrow a_1g_N=b_1g_D\Leftrightarrow bg_N=ag_D\Leftrightarrow M^*(C)$ is constant
over time. Neutral change is a proportional shift of both $A$ and $B$; a shift in $E$ is not technical change of this
kind (it is not a multiplicative shift of any homogeneous output index).
\item[(iv)] If $\kappa=1$ is imposed, then $S=\gamma/(ab)$ and every object in (ii), including $\mathcal B$, is identified, by
functional form.
\end{enumerate}
\end{proposition}

\begin{proof}
(i) At given compute the optimum solves $r(n,d,t)\equiv\ln(F_n/F_d)=0$ with $n+d=c-\ln6$. Differentiating,
$dn=-dd$ and $(r_d-r_n)\,dd+r_t\,dt=0$. At a point with $F_n=F_d$ the isoquant is tangent to the isocost, so the
derivative of $r$ in the direction $(-1,1)$ is the same along both curves. Along the isoquant,
$\sigma=d\ln M/d\ln\mathrm{MRTS}$ and $\ln\mathrm{MRTS}=\ln w+\ln M$ imply $d\ln w=(1/\sigma-1)\,d\ln M$; since
$d\ln M=2\,dd$, $r_d-r_n=2(1/\sigma^*-1)$. Hence $d\ln M^*/dt=2\,dd/dt=-r_t\sigma^*/(1-\sigma^*)$, and $r_t=\mathcal B_t$ because
$\mathrm{MRTS}=(F_n/F_d)(D/N)$. Lemma~\ref{lem:soc} gives $\sigma^*\in(0,1)$.
(ii) Apply Lemma~\ref{lem:alloc} with $\psi_t=g t$ to the inner aggregator: $\ln K_{1,t}$ falls at
$(a_1b_1/S)(g_N+g_D)$, so $\ln K_t=\kappa\ln K_{1,t}$ falls at $\gamma(g_N+g_D)$; $\ln G_t=\ln G_0+(b_1g_D-a_1g_N)t/S$. The
Hicks bias is $\partial_t\ln(a_1u/b_1v)$ at fixed $(n,d)$. Neither drift depends on $S$ once $(a,\gamma)$ are fixed, so
the two drifts give two linear equations in $(g_N,g_D)$ (for $g_N+g_D$ and $ag_D-bg_N$) with determinant $a+b=1$. A separate
Hicks-neutral trend $\omega_t=g_\omega t$ does not break this: in member $S$ it equals augmentation at
$(g_N,g_D)=(a,b)\,g_\omega/\gamma$, since $\kappa_Sa_1=\gamma/a$ and $\kappa_Sb_1=\gamma/b$. This is the same for every $S$, so
$(g_N,g_D)$ are the total augmentation rates, with the neutral part included. Observational equivalence follows
from Proposition~\ref{prop:fd}(iv) date by date. (iii) $M^*(C)=G_t^{-2}(C/6)^{b-a}$; $\mathcal B=0$ iff $a_1u/b_1v$ is constant
at fixed inputs, iff $A_t/B_t$ is constant. (iv) Immediate from $\gamma=abS$ when $\kappa=1$.
\end{proof}

\begin{remark}[Relation to Diamond, McFadden and Rodriguez]
\citet{diamond1978measurement} show that time series of output, inputs and marginal products cannot separate $\sigma$
from the bias of technical change. With constant relative input prices, the drift of the input ratio still reveals the
sign of the Hicks bias, because $0<\sigma^*<1$ (part (i)); in scaling laws the FLOP accounting fixes relative log-cost weights
at one, so the premise holds. With a CES technology the bias is $(1-1/\sigma)(g_N-g_D)$, and $\sigma^*<1$
(Corollary~\ref{cor:complements}) would then also reveal which input is augmented faster. That second step does not carry
over to the Chinchilla family, where $\mathcal B=\beta g_D-\alpha g_N$ (with $\kappa=1$): the drift signs
$\beta g_D-\alpha g_N$, which coincides with the sign of $g_D-g_N$ only when $\alpha=\beta$. For example, with
$(\alpha,\beta)=(0.30,0.45)$ and $(g_N,g_D)=(0.5,0.4)$, $\mathcal B=0.03>0$ and the compute-optimal $D/N$ falls, although
parameters are augmented faster. Within the family $g_N$ and $g_D$ are nonetheless identified separately (part (ii)), but
through the family's functional form, not through the drift alone. Proposition~\ref{prop:A-dmr} is a methodological
result; it is not applied to vintage data here, because published compute-optimal laws from different dates differ in
tokenizers, parameter conventions and outputs. The argument holds relative prices fixed. If the effective cost of tokens
relative to parameters moves over time, the drift of $M^*$ mixes price responses with bias, and the DMR problem returns. Such
moves come from data-acquisition costs, a tightening data constraint (Proposition~\ref{prop:A-wedge}(iii)(d)), or inference demand.
The sign result therefore needs training-optimal allocations (IsoFLOP minima) under FLOP-only costs at each date. A designed IsoFLOP sweep at each vintage identifies $\sigma^*_t$ and
hence $|\mathcal B_t|$ \citep[compare the insignificant, negatively correlated individual augmentation terms with a
well-identified sum in][]{ho2024algorithmic}.
\end{remark}

\begin{corollary}[Compute-equivalent gains; used in Online Appendix~E]\label{cor:ceg}
Compare an old and a new technology, each evaluated on its own frontier $L^*_j(C)=E_j+K_j(C/6)^{-\gamma_j}$. The
compute-equivalent gain $f(C)\equiv C^*_{\mathrm{old}}(L^*_{\mathrm{new}}(C))/C$ is constant in $C$ if and only if
$E_{\mathrm{old}}=E_{\mathrm{new}}$ and $\gamma_{\mathrm{old}}=\gamma_{\mathrm{new}}$, in which case
$f=(K_{\mathrm{old}}/K_{\mathrm{new}})^{1/\gamma}$. The individual exponents may differ: for example $(\alpha,\beta)=(0.3,0.4)$
and $(0.4,0.3)$ have the same $\gamma$. Factor-augmenting change gives $f=e^{\psi_N+\psi_D}$, and Hicks-neutral change gives
$f=e^{\Delta\omega/\gamma}$. If instead both technologies are evaluated at a common allocation rule $D=mN$, $f$ is constant
only if $E$ and the exponent sets $\{\alpha,\beta\}$ coincide, and, given that, if and only if the new law is the old one
with $N$ rescaled, $L_{\mathrm{new}}(N)=L_{\mathrm{old}}(\lambda N)$ for a fixed $\lambda$, in which case $f=\lambda^2$. Equal exponents
are therefore necessary but not sufficient at a common rule. Factor augmentation gives a constant $f$ on a common ray only if
$\psi_N=\psi_D$, and Hicks-neutral change does so only if $\alpha=\beta$.
\end{corollary}

\begin{proof}
With $x=C/6$, $\ln f=\gamma_o^{-1}[\ln K_o-\ln(E_n-E_o+K_nx^{-\gamma_n})]-\ln x$ and
$d\ln f/d\ln x=(\gamma_n/\gamma_o)K_nx^{-\gamma_n}/(E_n-E_o+K_nx^{-\gamma_n})-1$. This vanishes for all $x$ if and only if
$E_n=E_o$ and $\gamma_n=\gamma_o$. Factor augmentation multiplies $K$ by $e^{-\gamma(\psi_N+\psi_D)}$ (Lemma~\ref{lem:alloc}(iii)),
and $\omega$ multiplies it by $e^{-\Delta\omega}$. On a common ray, $L_j(N)=E_j+A_jN^{-\alpha_j}+B_jm^{-\beta_j}N^{-\beta_j}$.
Because $C=6mN^2$ on the ray, a constant gain $f$ requires $L_n(N)=L_o(\lambda N)$ for all $N$ with $\lambda=f^{1/2}$;
linear independence of power functions then forces equal $E$ and equal exponent sets, and matching the coefficient on each
power $N^{-e}$ requires the same $\lambda$ for both terms. With equal exponents and $A_n=A_oe^{-\alpha\psi_N}$,
$B_n=B_oe^{-\beta\psi_D}$ this gives $\lambda=e^{\psi_N}=e^{\psi_D}$. With $A_n=A_oe^{-\Delta\omega}$ and $B_n=B_oe^{-\Delta\omega}$ it
gives $\lambda^{\alpha}=\lambda^{\beta}=e^{\Delta\omega}$, which requires $\alpha=\beta$ when $\Delta\omega\neq0$.
\end{proof}

\begin{proposition}[Farrell decomposition]\label{prop:farrell}
Let run $i$ have compute-equivalent productivity $\Omega_i$ and true wedge $w_i$, and let $C_{\min,F}(L)$ be the compute
that best practice ($\Omega_F$) needs to reach loss $L$. Then
\[
\ln(6N_iD_i)-\ln C_{\min,F}(L_i)=\underbrace{\Phi(w_i)}_{\text{allocative}}+\underbrace{(\Omega_F-\Omega_i)/\gamma}_{\text{technical}}+\epsilon_i/\gamma,
\]
with $\Phi(w)=\ln(C/C_{\min})$ given by Lemma~\ref{lem:ce}. For two models at the same compute, the log compute-equivalent
gain is $\Delta\Omega/\gamma-\Delta\Phi$. In particular, a reallocation that keeps $\Omega$ fixed, such as the
Kaplan-to-Chinchilla move from Gopher-like to Chinchilla-like allocations, is purely allocative. Under the
\citet{besiroglu2024chinchilla} parameters, $\Phi(w_{\text{Gopher}})=\ln 2.01$.
\end{proposition}

\begin{proof}
For a lab with $(\omega,\psi)$, the minimal compute that reaches $L$ is $6[Ke^{-\Omega}/(L-E)]^{1/\gamma}$. This follows
from Lemma~\ref{lem:alloc}(iii): effective compute is $Ce^{\psi_N+\psi_D}$ and $R$ is scaled by $e^{-\omega}$. Hence
$\ln C_{\min,i}(L)-\ln C_{\min,F}(L)=(\Omega_F-\Omega_i)/\gamma$. Evaluating at the observed $L_i$, whose reducible part is
$e^{-\omega_i+\epsilon_i}R_i$, and using $C_i/C_{\min,i}(e^{-\omega_i}R_i)=e^{\Phi(w_i)}$ (Lemma~\ref{lem:ce} in
effective units) gives the decomposition. The Gopher number applies \eqref{eq:A-ce} at Gopher's wedge under those
parameters, $w=0.362$ ($N=280$B, $D=300$B).
\end{proof}

\subsection{Observational data: transmission, selection and proxies}\label{app:observational}

Throughout this subsection runs are on the expansion path of a common technology with Hicks-neutral productivity, so
$y_i=\omega_i+\gamma c_i-\epsilon_i$ up to a constant (Lemma~\ref{lem:alloc}(iii)). Compute follows
$c_i=\pi_0+\pi_1\omega_i+\eta_i$ with $(\omega,\eta,\epsilon)$ mutually uncorrelated, variances $V_\omega,V_\eta,V_\epsilon$.
Special cases are the \emph{target rule} (every lab trains to a common loss: $\pi_1=-1/\gamma$), \emph{exogenous budgets}
($\pi_1=0$) and \emph{funding} (more productive labs raise more compute: $\pi_1>0$).

\begin{proposition}[Transmission bias; used in Online Appendix~E]\label{prop:A-transmission}
Under the assumptions of this subsection:
\begin{enumerate}
\item[(i)] The OLS slope of $y$ on $c$ converges to $b_f=\gamma+\pi_1V_\omega/(\pi_1^2V_\omega+V_\eta)$, and its bias has the
sign of $\pi_1$. Under the target rule, $b_f=\gamma V_\eta/(V_\omega/\gamma^2+V_\eta)\in[0,\gamma)$. If all labs hit the
target exactly ($V_\eta=0$), $y$ is constant up to noise and cross-lab data carry no information about $\gamma$.
\item[(ii)] The $\omega$-component of the residual variance is $V_\omega V_\eta/\Var(c)<V_\omega$ whenever $\pi_1\neq0$:
TFP dispersion is understated.
\item[(iii)] The reverse-regression estimate $\gamma_r\equiv\Var(y)/\Cov(c,y)$ satisfies
\[
(\gamma_r-\gamma)\Cov(c,y)=(1+\gamma\pi_1)V_\omega+V_\epsilon.
\]
Hence, when $\Cov(c,y)>0$, the interval $[b_f,\gamma_r]$
contains $\gamma$ if and only if $-(1+V_\epsilon/V_\omega)/\gamma\le\pi_1\le0$. The bounds hold under target and
exogenous-budget behavior and fail under the funding rule, where both estimates exceed $\gamma$. The forward--reverse
bracket familiar from errors-in-variables problems \citep{klepper1984consistent} therefore does not survive
transmission that runs from productivity to compute.
\item[(iv)] The OLS slope of $n$ on $c$ equals $a$ for any $\pi_1$: Hicks-neutral transmission biases returns to compute,
not the allocation exponent. With factor-biased productivity, the slope converges to
$a+\Cov(\chi,c)/[(\alpha+\beta)\Var(c)]$.
\item[(v)] (Panel remedies.) With lab fixed effects and time-invariant $\omega_f$, the within estimator is consistent
\citep{mundlak1961empirical}. With $\omega_{ft}=\rho_\omega\,\omega_{f,t-1}+\xi_{ft}$ and compute chosen before $\xi_{ft}$ is
realized, $y_{ft}=\lambda_0(1-\rho_\omega)+\gamma c_{ft}+\rho_\omega(y_{f,t-1}-\gamma c_{f,t-1})+\xi_{ft}-\epsilon_{ft}+\rho_\omega\epsilon_{f,t-1}$.
The moments $\E[(\xi_{ft}-\epsilon_{ft}+\rho_\omega\epsilon_{f,t-1})z]=0$ with $z\in\{1,c_{ft},c_{f,t-1},y_{f,t-2}\}$ over-identify
$(\gamma,\rho_\omega,\lambda_0)$, where $\lambda_0$ is the intercept of $y$. The instrument $c_{ft}$ is valid only if compute does not
respond to the previous run's noise $\epsilon_{f,t-1}$ (as when the lab knows $\omega_{f,t-1}$ and treats $\epsilon$ as
evaluation noise). If the budget responds to the realized outcome $y_{f,t-1}$, drop $c_{ft}$ and use
$\{1,c_{f,t-1},c_{f,t-2},y_{f,t-2}\}$. The first-stage proxy inversion of \citet{olley1996dynamics} and
\citet{levinsohn2003estimating} is not available here (Proposition~\ref{prop:proxy}), so the lagged-outcome form replaces it.
\end{enumerate}
\end{proposition}

\begin{proof}
(i) $\Cov(y,c)=\gamma\Var(c)+\pi_1V_\omega$ and $\Var(c)=\pi_1^2V_\omega+V_\eta$; substitute $\pi_1=-1/\gamma$ for the
target rule. With $V_\eta=0$ under the target rule, $y=\text{const}-\epsilon$. (ii) The residual variance is
$\Var(y)-\Cov(y,c)^2/\Var(c)=V_\omega+V_\epsilon-\pi_1^2V_\omega^2/\Var(c)$. (iii) $\Var(y)-\gamma\Cov(c,y)=\gamma\pi_1V_\omega+V_\omega+V_\epsilon$.
$b_f\le\gamma$ iff $\pi_1\le0$; $\gamma_r\ge\gamma$ iff $(1+\gamma\pi_1)V_\omega+V_\epsilon\ge0$. (iv) By
Lemma~\ref{lem:alloc}(iii), $n=\ln G+a(c-\ln6)+\chi/(\alpha+\beta)$ exactly, and $\omega$ does not enter. (v) Substitute
$\omega_{f,t-1}=y_{f,t-1}-\gamma c_{f,t-1}+\epsilon_{f,t-1}$ (up to the constant). Here $c_{ft}$ and $c_{f,t-1}$ are chosen
before $\xi_{ft}$ and $\epsilon_{ft}$ are realized. $c_{f,t-1}$ is also chosen before $\epsilon_{f,t-1}$ is realized, while $c_{ft}$ is
uncorrelated with $\epsilon_{f,t-1}$ under the stated behavioral condition. $y_{f,t-2}$ is uncorrelated with
$(\xi_{ft},\epsilon_{ft},\epsilon_{f,t-1})$ when $\epsilon$ is serially independent.
\end{proof}

\begin{proposition}[Selection on the outcome; used in Online Appendix~E]\label{prop:selection}
Suppose a run is released (reported) only if $y_i\ge\bar y$. Let $S_i\equiv\omega_i-\epsilon_i$ and
$m(k)\equiv\E[S\mid S\ge k]$.
\begin{enumerate}
\item[(i)] (Exogenous compute, $\pi_1=0$.) In the released sample $\E[y\mid c]=\gamma c+m(\bar y-\gamma c)$, whose slope is
$\gamma[1-m'(\bar y-\gamma c)]$. Since $m'\ge0$ for every distribution, the released-sample OLS slope is at most $\gamma$.
If $S$ has a log-concave density, then $0\le m'\le1$ \citep{bagnoli2005logconcave} and the slope lies in $[0,\gamma]$;
without log-concavity it can be negative. If moreover $\E[\omega\mid S]$ is nondecreasing (for example $\omega$ and
$\epsilon$ independent with log-concave densities, \citealp{efron1965increasing}), then $\E[\omega\mid c,\text{released}]$
is nonincreasing in $c$.
\item[(ii)] (Gaussian case, any $\pi_1$.) Let $(\omega,\eta,\epsilon)$ be jointly normal, $\delta\equiv\Cov(c,y)/\Var(y)$ and
$V_{c|y}\equiv\Var(c)-\delta^2\Var(y)$. Selection on $y$ leaves the reverse regression of $c$ on $y$ unchanged
($\delta_{\text{sel}}=\delta$). The forward slope becomes
\[
b_{\text{sel}}=\frac{\delta\,r\,\Var(y)}{\delta^2r\Var(y)+V_{c|y}},
\qquad r\equiv\frac{\Var(y\mid y\ge\bar y)}{\Var(y)},
\]
where $r=1+z\lambda(z)-\lambda(z)^2\in(0,1)$
with $z=(\bar y-\E y)/\mathrm{sd}(y)$ and $\lambda(z)=\phi(z)/[1-\Phi(z)]$. The slope $b_{\text{sel}}$ is increasing in $r$.
\item[(iii)] (Net sign.) If $\pi_1\le0$, then $b_{\text{sel}}\le\gamma$, and $b_{\text{sel}}\le b_f\le\gamma$ when
$\Cov(c,y)>0$: the net bias is negative. If $\pi_1>0$, then
$\gamma\delta<1$ and $b_{\text{sel}}>\gamma$ if and only if
\[
r>r^*\equiv\frac{\gamma\,V_{c|y}}{\delta\Var(y)(1-\gamma\delta)}.
\]
Mild selection leaves the upward transmission bias dominant; severe selection reverses it.
\item[(iv)] Since selection on $y$ lowers the forward slope and leaves $1/\delta$ unchanged, the bracket of
Proposition~\ref{prop:A-transmission}(iii) remains valid in released samples whenever
$-(1+V_\epsilon/V_\omega)/\gamma\le\pi_1\le0$.
\end{enumerate}
\end{proposition}

\begin{proof}
(i) Given $c$ (independent of $S$), released iff $S\ge\bar y-\gamma c$. $m$ is nondecreasing because
$\E[S\mid S\ge k_1]$ is a mixture of $\E[S\mid k_1\le S<k_2]\le k_2$ and $\E[S\mid S\ge k_2]\ge k_2$ for $k_1<k_2$. The CEF
is therefore $\gamma c$ plus a nonincreasing function of $c$, and the linear projection of a nonincreasing function on $c$
has a nonpositive slope. Log-concavity implies a nonincreasing mean residual life $m(k)-k$, hence $m'\le1$. For the last
claim, $\E[\omega\mid S\ge k]=\E[h(S)\mid S\ge k]$ with $h$ nondecreasing, and $S\mid S\ge k$ is stochastically increasing in
$k=\bar y-\gamma c$. (ii) Under joint normality, $c=\E c+\delta(y-\E y)+\nu$ with $\nu$ independent of $y$, so selection
on $y$ does not change the distribution of $\nu$ or its independence from $y$. Then $\Cov_{\text{sel}}(c,y)=\delta\Var_{\text{sel}}(y)$
and $\Var_{\text{sel}}(c)=\delta^2\Var_{\text{sel}}(y)+V_{c|y}$. The truncated-normal variance gives $r$. Differentiating in
$r$ gives a derivative with the sign of $\delta V_{c|y}>0$. (iii) For $\pi_1\le0$ and $\delta>0$ combine (ii) with Proposition~\ref{prop:A-transmission}(i); if $\delta\le0$,
then $b_{\text{sel}}\le0<\gamma$ directly. For $\pi_1>0$, $\Var(y)-\gamma\Cov(c,y)=(1+\gamma\pi_1)V_\omega+V_\epsilon>0$ gives
$\gamma\delta<1$; rearranging $b_{\text{sel}}>\gamma$ gives the threshold. (iv) follows from (ii) and
Proposition~\ref{prop:A-transmission}(iii).
\end{proof}

\begin{remark}
Part (i) is the classical attenuation of least squares under truncation on the dependent variable
\citep{goldberger1981linear,hausman1977social}; the new elements are its interaction with transmission (parts (ii)--(iii))
and the survival of the forward--reverse bracket (part (iv)). The model specification's statement that
$\E[\omega\mid c,\text{released}]$ is decreasing in $c$ holds for exogenous compute (part (i)). Under the funding rule, $\E[\omega\mid c,\text{released}]=\E[\omega\mid c]+\Delta(c)$, where only the
selection component $\Delta$ is decreasing; with mild selection the total is increasing (verified by simulation).
Dropping the highest-loss runs from a designed sweep \citep{besiroglu2024chinchilla} is truncation on the dependent
variable of the same kind, with $\pi_1=0$ and selection operating through $\epsilon$ and the design.
\end{remark}

\begin{proposition}[Allocation-based proxies; used in Online Appendix~E]\label{prop:proxy}
Suppose runs minimize cost with a common and known wedge $\bar w$ (for example $T=0$), share $(A,B,\alpha,\beta)$, and
differ in unobserved $(\omega_i,\psi_{N,i},\psi_{D,i})$.
\begin{enumerate}
\item[(i)] $\ln M_i=m_0(c_i)+\frac{2}{\alpha+\beta}(\ln\bar w-\chi_i)$ with $m_0(c)=-2\ln G+(b-a)(c-\ln6)$. The pair $(c,M)$
reveals the bias index exactly, $\chi_i=\ln\bar w-\tfrac{\alpha+\beta}{2}[\ln M_i-m_0(c_i)]$: a
Levinsohn--Petrin-type invertibility holds with a known linear inverse.
\item[(ii)] Given compute, $M$ does not depend on $\omega$ or on any productivity with $\chi=0$. No allocation-based
proxy exists for Hicks-neutral productivity: the policy function is flat in $\omega$.
\item[(iii)] Even for $\chi$, the proxy step does not identify the technology. The map $(n,d)\mapsto(c,\ln M)$ is a
bijection, so any first-stage control function $\varphi(c,\ln M)$ absorbs $F(n,d)$ entirely; this is the
functional-dependence problem of \citet{ackerberg2015identification}. Identification must come from timing
restrictions (for example, $\chi$ Markov and compute chosen before the innovation) or from a known mapping between $\chi$
and level productivity $\Omega$. The same holds for compute as a proxy for $\omega$ under a strictly monotone budget rule
$c=h(\omega,z)$, because on the path the inputs are functions of $c$ alone.
\item[(iv)] If wedges are heterogeneous, the allocation reveals only $\ln w_i-\chi_i$. Inference demand and factor bias
cannot be separated from $(N,D)$ alone; this is the contamination in Proposition~\ref{prop:A-wedge}(ii).
\end{enumerate}
\end{proposition}

\begin{proof}
(i)--(ii) Lemma~\ref{lem:geometry}(iii) and Proposition~\ref{prop:A-wedge}(ii): the econometrician's
$\hat w=(M/M^*_0(C))^{(\alpha+\beta)/2}$, computed with $\chi=0$, equals $we^{-\chi}$; solve for $\ln M$. The derivative in
$\omega$ is zero by Lemma~\ref{lem:alloc}(iv). (iii) $n=\tfrac12(c-\ln6-\ln M)$ and $d=\tfrac12(c-\ln6+\ln M)$. (iv) is (i)
with $\bar w$ replaced by $w_i$.
\end{proof}

\subsection{What over-training reveals: the generalized wedge}\label{app:wedge}

\begin{proposition}[Generalized wedge; Proposition~\ref{prop:wedge} in the main text]\label{prop:A-wedge}
A developer chooses $(N,D)$ to maximize
\begin{equation}\label{eq:A-gw-obj}
\begin{gathered}
\Pi(N,D)=V\big(L(N,D),N\big)-X_T(N,D)-c_DD\\
\textup{subject to}\quad g_k(n,d)\le0,\quad k=1,\dots,K,
\end{gathered}
\end{equation}
where $V$ is differentiable with $\partial V/\partial L<0$; $X_T\equiv c_T(N)\,6ND$ is training expenditure at the
(shadow) price $c_T(N)$ of a training FLOP, whose size elasticity is $\delta\equiv d\ln c_T/d\ln N$; $c_D\ge0$ is a cost
per training token; $L$ is given by \eqref{eq:A-tech} or by any technology with $\varepsilon_N,\varepsilon_D>0$; and the
$g_k$ are differentiable constraints. Let $m_N\equiv-(\partial V/\partial\ln N)_L/X_T$ be the value of compactness at
fixed quality per unit of training expenditure, $m_D\equiv c_DD/X_T$, and $\xi_k\ge0$ the multipliers (value per log
point).
\begin{enumerate}
\item[(i)] At any interior Karush--Kuhn--Tucker point,
\begin{equation}\label{eq:A-gw}
w=\frac{1+\delta+m_N+\sum_k\xi_kg_{k,n}/X_T}{1+m_D+\sum_k\xi_kg_{k,d}/X_T}.
\end{equation}
With $\delta=m_D=0$ and no binding constraint, $w-1=m_N$, where $m_N=-(\partial V/\partial\ln N)_L/(c_T\,6ND)$
is the marginal value of compactness at fixed quality divided by training expenditure. The expenditure share
$s\equiv(w-1)/w$ is then $m_N/(1+m_N)$.
\item[(ii)] (Invariance and contamination.) $w$ is computed from $(N,D)$ and the technology alone; it is invariant to
$\omega$, $\epsilon$, $E$ and $\kappa$. If the developer's first-order conditions involve a technology $L'$ (its own
productivity, a law it believes, or its own output concept) while the econometrician uses $L$, then
$\hat w=w\cdot w_L(N,D)/w_{L'}(N,D)$. With factor-augmenting productivity $w_L/w_{L'}=e^{-\chi}$, so $\hat w=we^{-\chi}$.
\item[(iii)] (Special cases.) Let $T$ be lifetime inference tokens in present value ($T=q/(r+h)$ for a token flow $q$,
discount rate $r$ and obsolescence hazard $h$), $c_S(N)$ the price of a serving FLOP, $X_S\equiv c_S\,2NT$ serving
expenditure and $p\equiv c_S/c_T$ the price ratio.
\begin{enumerate}
\item[(a)] (Serving cost borne at rate $\lambda$.) If $V=R(L)-\lambda X_S$ with $T=T(L,N)$, then
$w=1+\delta+\lambda p(\phi+\eta)\,T/(3D)$, where $\phi\equiv1+d\ln c_S/d\ln N$ and $\eta\equiv(\partial\ln T/\partial\ln N)_L$;
a direct effect of size on revenue at fixed quality adds $-(\partial R/\partial\ln N)_L/X_T$. With $\delta=\eta=0$ and
$\phi=1$: $w=1+\lambda pT/(3D)$, $w-1=\lambda X_S/X_T$ and $s=\lambda X_S/(X_T+\lambda X_S)$, the internalized serving share
of lifetime expenditure; no conversion of FLOPs into dollars is needed. $\lambda=1$ for a developer that serves its own
model, $\lambda<1$ when users bear part of the cost. Compute proportional to $N$ that is planned at training time
(reinforcement-learning rollouts, synthetic-data generation, evaluation) enters $X_S$ in the same way.
\item[(a$'$)] (Adoption.) If the developer bears no serving cost and $V=v(A)$, where adoption $A(L,P_u)$ falls with users'
serving cost per token $P_u\propto N$, then $m_N=v'(A)A\,|\varepsilon_{A,P}|/X_T$: $w-1$ is the adoption value of
compactness, not a token count. It is case (a) with $\lambda\equiv v'(A)A|\varepsilon_{A,P}|/X_u$, where $X_u$ is users'
serving expenditure.
\item[(b)] (Tier or memory cap.) If $N\le\bar N$ binds, $w=(1+\nu)(1+\delta+m_N)/(1+m_D)$ with
$\nu\equiv\xi/[X_T(1+\delta+m_N)]\ge0$. Any $w\ge(1+\delta+m_N)/(1+m_D)$ is rationalizable, including with $m_N=0$, so
$w-1$ bounds $m_N$ from above (when $\delta=m_D=0$) but does not identify it.
\item[(c)] (Latency.) A binding service-level objective $\tau(N)\le\bar\tau$, with elasticity
$\ell\equiv d\ln\tau/d\ln N>0$, is case (b) with $\nu=\ell\xi_\tau/[X_T(1+\delta+m_N)]$. A latency cost inside $V$
instead enters $m_N$.
\item[(d)] (Data cost and data cap.) $w=(1+\delta+m_N)/(1+m_D)$, so a data cost lowers $w$ for given $m_N$. With
$\delta=0$ and $\lambda=p=1$, $\hat T/D\equiv3(w-1)=(T/D-3m_D)/(1+m_D)<T/D$, and $s=(X_S-X_D)/(X_T+X_S)$ with
$X_D\equiv c_DD$. A binding cap $D\le\bar D$ adds $\mu\equiv\xi_D/X_T\ge0$ to the denominator,
$w=(1+\delta+m_N)/(1+m_D+\mu)$, which can fall below one with $T=0$.
\item[(e)] (Logit distillation.) Teacher forward passes cost $2N_TD$ FLOPs at the training price, a per-token cost with
$m_D=N_T/(3N)$: with $\delta=0$ and $\phi=1$, $w=\big(1+\lambda pT/(3D)\big)/\big(1+N_T/(3N)\big)$.
\item[(f)] (Utilization varying with size.) If the prices of training and serving FLOPs have size elasticities $\delta$
and $\delta_S$ ($\phi=1+\delta_S$), then $w=1+\delta+\lambda(1+\delta_S)pT/(3D)$, which equals
$(1+\delta)\big(1+\lambda pT/(3D)\big)$ if and only if $\delta_S=\delta$.
\end{enumerate}
\item[(iv)] (Signs.) Let $w_0\equiv(1+\delta+m_N)/(1+m_D)$. A single binding constraint raises $w$ above $w_0$ if and only
if $g_n>w_0\,g_d$. Memory or tier caps and latency objectives raise $w$; data caps and
deadlines that bind on tokens lower it; deadlines proportional to compute ($g_n=g_d$) pull it toward one; minimum-size
floors ($g_n<0=g_d$) lower it. Hence $\hat T\equiv3D(\hat w-1)$ overstates $\lambda pT$ under caps, latency objectives,
$\delta>0$ (with $\delta_S\ge0$) and $\chi<0$, and understates it under data costs, data caps, distillation, minimum-size
floors, $\delta<0$ (with $\delta_S\le0$) and $\chi>0$. A wedge below one cannot be generated by serving demand as long as
serving expenditure rises with size at given quality ($\phi+\eta\ge0$).
\end{enumerate}
\end{proposition}

\begin{proof}
In $(n,d)$ the Lagrangian is $\Pi-\sum_k\xi_kg_k$. Since $\partial L/\partial n=-(L-E)\varepsilon_N$ and
$\partial L/\partial d=-(L-E)\varepsilon_D$, with $\Lambda\equiv-(\partial V/\partial L)(L-E)>0$ and
$\partial X_T/\partial n=(1+\delta)X_T$, $\partial X_T/\partial d=X_T$, the first-order conditions are
$\Lambda\varepsilon_N+(\partial V/\partial\ln N)_L-(1+\delta)X_T-\sum_k\xi_kg_{k,n}=0$ and
$\Lambda\varepsilon_D-X_T-c_DD-\sum_k\xi_kg_{k,d}=0$. Dividing the first by the second gives \eqref{eq:A-gw}.
(ii) Remark~\ref{rem:exact}; in the generalized family $w=a_1u/(b_1v)$ does not involve $\kappa$. If the first-order
conditions involve $L'$, \eqref{eq:A-gw} determines $w_{L'}(N,D)$, and the econometrician computes
$w_L(N,D)=w_{L'}(N,D)\cdot w_L/w_{L'}$. With factor-augmenting productivity the true elasticities use
$u=Ae^{-\alpha(n+\psi_N)}$ and $v=Be^{-\beta(d+\psi_D)}$, so $w_{L'}/w_L=e^{\beta\psi_D-\alpha\psi_N}=e^{\chi}$.
(iii)(a) $-(\partial V/\partial\ln N)_L=\lambda\,\partial X_S/\partial\ln N|_L=\lambda X_S(\phi+\eta)$ because
$\ln X_S=\ln c_S+\ln2+\ln N+\ln T$; divide by $X_T$ and use $X_S/X_T=pT/(3D)$. The share follows from $w-1=\lambda X_S/X_T$.
(a$'$) $(\partial V/\partial\ln N)_L=v'(A)\,\partial A/\partial\ln P_u$. (b) $g=n-\bar n$, so $g_n=1$ and $g_d=0$.
(c) $g=\ln\tau(n)-\ln\bar\tau$, so $g_n=\ell$. (d) Definition of $m_D$; the cap is $g=d-\bar d$. (e) The teacher's cost
$2N_TDc_T$ has elasticity zero in $N$ and one in $D$, so it acts as $c_D=2N_Tc_T$ and $m_D=2N_TD/(6ND)$. (f) The elasticity
of $X_T$ in $N$ is $1+\delta$ and that of $X_S$ is $1+\delta_S$, so $w=[(1+\delta)X_T+\lambda(1+\delta_S)X_S]/X_T$, and
$(1+\delta)(1+\lambda pT/(3D))-w=(\delta-\delta_S)\lambda pT/(3D)$, which vanishes iff $\delta_S=\delta$.
(iv) With $A\equiv1+\delta+m_N$ and $B\equiv1+m_D$, $\partial_\xi[(A+\xi g_n/X_T)/(B+\xi g_d/X_T)]$ has the sign of
$g_nB-g_dA$, that is, of $g_n-w_0g_d$, for every $\xi$. The first-order conditions are
necessary, and the inversion uses nothing else; they are also sufficient when $\Pi$ is concave in $(n,d)$, for example in
the lifetime-compute problem, where $\ln(6ND+2NT)$ is convex in logs and $\{y\ge\bar y\}$ is convex in the generalized
family.
\end{proof}

\begin{remark}[Conduct, expenditure shares, and the parallel with markup estimation]\label{rem:A-conduct}
Proposition~\ref{prop:A-wedge} inverts the first-order conditions of a profit-maximizing developer in the spirit of
\citet{hao2026theory}, with an imperfectly substitutable technology; in their Leontief limit over-training does not arise.
Cases (a), (a$'$) and (b) are three models of conduct: a developer that serves its own model and bears the serving cost
($\lambda=1$); an open-weight developer that internalizes users' costs only through adoption; and size-class competition
with $N$ set by a menu of deployment tiers. They make different predictions. Under (a), $w-1$ rises with the developer's
own serving footprint; under (a$'$), with the sensitivity of adoption to users' serving cost; under (b), sizes bunch at
tier boundaries and $w-1$ is unrelated to demand given the tier. The tools of \citet{berry2014identification},
\citet{backus2021common} and \citet{duarte2024testing} for testing conduct apply. Without knowing $p$, the object that (a)
identifies is the ratio of internalized serving expenditure to training expenditure, $w-1=\lambda\phi X_S/X_T$, or the
share $s$; a count of tokens needs $p$: $T=3D(w-1)/(\lambda\phi p)$.

The wedge is the ratio of the two inputs' output elasticities to their cost elasticities, which is the statistic
\citet{raval2023testing} compares across inputs, here with FLOP-accounting weights. Part (ii) contains the two classic
critiques of the production approach to markups: factor-augmenting heterogeneity \citep{demirer2020production,raval2023testing}
and a mismatch between the output the econometrician measures and the output the firm values
\citep{bond2020unpleasant}. \citet{doraszelski2021reexamining} show that consistent estimation of the output elasticity by
proxy methods generally requires controlling for the markup itself. Here the elasticities come from designed experiments,
not from proxy estimation on the same choices, which removes that circularity at the cost of the experimental
technology's external validity for each developer. As in \citet{deridder2026hitchhiker}, levels are fragile and ranks are
more robust: at given compute, the rank of $\hat w$ across models is the rank of $M$ for every member of the generalized
family.
\end{remark}

\begin{corollary}[Sufficient statistics; used in Proposition~\ref{prop:wedge}]\label{cor:suff}
In the generalized family, for a model with compute $C$ and ratio $M$ at an interior optimum of
Proposition~\ref{prop:A-wedge} with $\delta=m_D=0$ and no binding constraint,
\begin{gather*}
m_N=w-1=\Big(\frac{M}{M^*(C)}\Big)^{1/\sigma^*-1}-1,\qquad s=1-\Big(\frac{M}{M^*(C)}\Big)^{-(1/\sigma^*-1)},\\
\frac{C}{C_{\min}}=\Big(\frac{\alpha+\beta w}{\alpha+\beta}\Big)^{1/\gamma}w^{-1/\alpha},
\end{gather*}
where $M^*(C)$ is the training-optimal ratio at the model's own compute (the last formula with $\kappa=1$; in general
Lemma~\ref{lem:ce}). Under Proposition~\ref{prop:A-wedge}(iii)(a) with $\phi=1$ and $\eta=0$,
$\lambda pT/D=3[(M/M^*(C))^{1/\sigma^*-1}-1]$. In the homothetic case $\alpha=\beta=\rho$, $M^*=(B/A)^{1/\rho}$ is constant,
$1/\sigma^*-1=\rho$ and $C/C_{\min}=\cosh(\tfrac\rho2\ln(M/M^*))^{2/\rho}$. With $(\alpha,\beta)$ unequal, $\rho$ is
replaced by the mean exponent $(\alpha+\beta)/2=1/\sigma^*-1$ and $M^*$ by $M^*(C)$. This justifies evaluating a
technology's own compute-optimal law at each model's own compute.
\end{corollary}

\begin{proof}
Combine Proposition~\ref{prop:A-wedge}(i), equation~\eqref{eq:A-wM} and Lemma~\ref{lem:ce}.
\end{proof}

\begin{proposition}[Family token budgets; Proposition~\ref{prop:wedge} in the main text]\label{prop:A-family}
A developer trains $K$ models of sizes $N_1,\dots,N_K$ on a common token budget $D$ and chooses $D$ (and, where stated,
the sizes) to maximize $\sum_iV_i(L(N_i,D),N_i)-\sum_iX_{T,i}-c_DD$ subject to $n_i\le\bar n_i$, where
$X_{T,i}=c_T\,6N_iD$ and the data cost is incurred once. Let $w_i$ be member $i$'s wedge computed from the technology,
$\omega_i\equiv N_i/\sum_jN_j$ its compute share, $s_i\equiv(w_i-1)/w_i$, $m_{N,i}\equiv-(\partial V_i/\partial\ln N_i)_L/X_{T,i}$,
$m_D\equiv c_DD/\sum_jX_{T,j}$, $\Lambda_i\equiv-(\partial V_i/\partial L)(L_i-E)$, and
$\rho_i\equiv\Lambda_i\varepsilon_{D,i}/X_{T,i}$, the marginal value of data to member $i$ per unit of its marginal training
cost. Let $\bar w_H\equiv(\sum_i\omega_i/w_i)^{-1}$ be the compute-weighted harmonic mean of the member wedges and
$\pi_i\equiv(\omega_i/w_i)/\sum_j(\omega_j/w_j)$.
\begin{enumerate}
\item[(i)] (Member wedges are technological.) Given $D$, $\ln w_i-\ln w_j=\ln w(n_i,d)-\ln w(n_j,d)$, which equals
$-a_1(n_i-n_j)$ in the generalized family: the profile of $\hat T_i\equiv3D(w_i-1)$ across sizes is fixed by the
technology.
\item[(ii)] (Family first-order condition.) At an interior optimum in $D$, $\sum_i\omega_i\rho_i=1+m_D$. If the sizes are
set by a menu, this is the only first-order condition, and the $w_i$ carry no information about the $m_{N,i}$.
\item[(iii)] (Sizes chosen, caps slack.) If each $n_i$ is also chosen at an interior optimum,
$\sum_i\pi_i(1+m_{N,i})=(1+m_D)\,\bar w_H$, where $\bar w_H$ is computed from the technology alone and satisfies the
identity $1-1/\bar w_H=\sum_i\omega_is_i$ (so the family's first-order condition pins the $\pi$-weighted average of the
members' values of compactness, not $\bar w_H$, which is observed). Under
Proposition~\ref{prop:A-wedge}(iii)(a) with common $(\lambda,p)$, $\delta=\eta=0$, $\phi=1$ and $m_D=0$:
$\sum_i\pi_i\hat T_i=\lambda p\sum_i\pi_iT_i$. Only this $\pi$-weighted average of the members' inversions is a
revealed-preference object; the member-level $T_i$ are not identified (one equation for $K$ unknowns).
\item[(iv)] (Binding caps and floors.) If every size is at a binding cap, the equality in (iii) becomes
$\sum_i\pi_i(1+m_{N,i})\le(1+m_D)\bar w_H$, so that $\sum_i\pi_i\hat T_i\ge\lambda p\sum_i\pi_iT_i$: $\bar w_H-1$ bounds the
family's average value of compactness from above. Binding minimum-size floors reverse the inequality.
\item[(v)] (A common $D$ under member-level choice is a knife-edge.) If instead each member chooses its own $D_i$
(Proposition~\ref{prop:A-wedge} with $\delta=m_D=0$ and no binding constraint), equal token counts $D_i=\bar D$ occur only
if $1+m_{N,i}=w(n_i,\bar d)$ for every $i$; in the generalized family, only if
$(1+m_{N,i})/(1+m_{N,j})=(N_i/N_j)^{-a_1}$ exactly.
\end{enumerate}
\end{proposition}

\begin{proof}
(i) In the generalized family $\ln w=\ln(a_1A/b_1B)-a_1n+b_1d$ (Lemma~\ref{lem:geometry}(iii)). (ii) The first-order
condition in $d$ is $\sum_i\Lambda_i\varepsilon_{D,i}=\sum_iX_{T,i}+c_DD$; divide by $\sum_jX_{T,j}$ and note
$X_{T,i}/\sum_jX_{T,j}=\omega_i$. (iii) The condition in $n_i$ is $\Lambda_i\varepsilon_{N,i}=X_{T,i}(1+m_{N,i})$, so
$\rho_i=(1+m_{N,i})/w_i$; substitute into (ii) and multiply by $\bar w_H$. For the shares,
$1-\sum_i\omega_i/w_i=\sum_i\omega_i(1-1/w_i)$. With $m_D=0$, $\sum_i\pi_im_{N,i}=\bar w_H-1=\sum_i\pi_i(w_i-1)$ because
$\sum_i\pi_iw_i=\bar w_H$; multiply by $3D$ and use $m_{N,i}=\lambda pT_i/(3D)$. (iv) With a binding cap the condition in
$n_i$ is $\Lambda_i\varepsilon_{N,i}=X_{T,i}(1+m_{N,i})+\xi_i$ with $\xi_i\ge0$, so $\rho_i\ge(1+m_{N,i})/w_i$; a floor has
$\xi_i\le0$. (v) Proposition~\ref{prop:A-wedge}(i) member by member, and (i).
\end{proof}

\begin{remark}
In a brute-force example with three members whose sizes are chosen on a common budget, the member inversions
$\hat T_i/T_i$ are 0.86, 0.91 and 1.06 while the $\pi$-weighted average is exact (replication package). Families that train
every size on the same budget are therefore informative about serving demand only at the family level, and only if the
sizes are chosen at the margin; if sizes come from a menu of tiers they are informative about nothing beyond the bound in
(iv).
\end{remark}

\subsection{Model-free identification of the elasticity of substitution and the wedge}\label{app:modelfree}

\begin{proposition}[Model-free $\sigma^*$ and $w$; Proposition~\ref{prop:modelfree} in the main text]\label{prop:A-modelfree}
Let $y=F(n,d)$ be $C^2$ with $F_n,F_d>0$, where $y$ is any strictly increasing transformation of $-L$. Write
$x\equiv\ln M=d-n$ and $c\equiv\ln C=\ln6+n+d$, let $Y(x,c)\equiv F\big((c-\ln6-x)/2,\,(c-\ln6+x)/2\big)$, let $L$ also
denote loss as a function of $(x,c)$, and let $L^*(c)\equiv\min_xL(x,c)$ be the loss--compute frontier with argmin
$x^*(c)=\ln M^*(C)$.
\begin{enumerate}
\item[(i)] (Wedge.) At any point, $F_n=Y_c-Y_x$, $F_d=Y_c+Y_x$ and
\[
w=\frac{Y_c-Y_x}{Y_c+Y_x}=\frac{L_c-L_x}{L_c+L_x}.
\]
$w$ is invariant to strictly monotone transformations of output, so it requires neither $E$ nor $\kappa$ nor a
functional form: only the slopes of any output measure along the isocost and along the ray of constant $M$.
\item[(ii)] (Elasticity of substitution at the compute optimum.) At $x=x^*(c)$, where $L_x=0$,
\begin{equation}\label{eq:A-mf}
\frac1{\sigma^*}-1=\frac{2(-Y_{xx})}{Y_c}=\frac{2L_{xx}}{|dL^*/dc|}=\frac{L_{nn}|_C}{2\,|dL^*/dc|},
\end{equation}
where $L_{nn}|_C$ is the second derivative of the IsoFLOP profile with respect to $\ln N$ at its minimum and
$dL^*/dc=L_c$ by the envelope theorem. The ratio is invariant to strictly monotone transformations of output at the
argmin, but not elsewhere: the curvature of the IsoFLOP profile divided by twice the slope of the loss--compute frontier,
in any common output units, identifies $\sigma^*$.
\item[(iii)] (Slope of the log wedge.) At $x^*(c)$, $\partial\ln w/\partial\ln M|_C=1/\sigma^*-1$. In the generalized
family $\ln w$ is exactly linear in $\ln(M/M^*(C))$ with this slope at every compute (Lemma~\ref{lem:geometry}(iii)); in
general the slope varies with $M$, which gives a test of the family.
\item[(iv)] (Local $\sigma$ anywhere.) $F_{nn}=Y_{cc}-2Y_{cx}+Y_{xx}$, $F_{dd}=Y_{cc}+2Y_{cx}+Y_{xx}$ and
$F_{nd}=Y_{cc}-Y_{xx}$, so $\sigma=P/(P+Q)$ (Lemma~\ref{lem:soc}) is identified at any point from the second-order expansion
of any output in $(x,c)$.
\item[(v)] (What $\kappa=1$ imposes.) In the generalized family, with $y=-\ln(L-E)$, $dy^*/dc=\gamma$ and
$-y_{nn}|_C=S\gamma$ at the argmin, so \eqref{eq:A-mf} returns $S/2$ for every $\kappa$ and $E$. The restriction $\kappa=1$
imposes $S=\gamma/(ab)$, that is, $-y_{nn}|_C=\gamma^2/(ab)$: the curvature of every IsoFLOP profile in log reducible loss is
pinned by the slopes of the frontier and of the path. With $\kappa$ free the curvature is a separate parameter, which
\eqref{eq:A-mf} measures without the restriction.
\item[(vi)] (Designs.) $\sigma^*$ is identified by IsoFLOP profiles: the curvature of a profile at its minimum, and the
frontier slope from the minima at two or more budgets (or locally). On-path designs identify the frontier slope but not
the curvature, hence not $\sigma^*$ without a functional form (Proposition~\ref{prop:fd}). The wedge at a point is identified by local variation in two non-collinear directions around
it (for example, a local factorial in $N$ and $D$), and only inside the design's support; its precision falls as
$\varepsilon_D\to0$, since $w\to\infty$ when $L_c+L_x\to0$.
\end{enumerate}
\end{proposition}

\begin{proof}
(i) $\partial/\partial n=\partial_c-\partial_x$ and $\partial/\partial d=\partial_c+\partial_x$. With $L=h(Y)$ and $h'<0$,
$L_x=h'Y_x$ and $L_c=h'Y_c$, and $h'$ cancels; any strictly monotone $g$ multiplies both first derivatives by $g'$.
(ii) At $F_n=F_d\equiv f$, $P=2f^3$ and $Q=-f^2(F_{nn}-2F_{nd}+F_{dd})=-4f^2Y_{xx}$, since
$Y_{xx}=(F_{nn}-2F_{nd}+F_{dd})/4$. By Lemma~\ref{lem:soc}, $1/\sigma-1=Q/P=-2Y_{xx}/f$ with $f=Y_c$. For $L=h(Y)$,
$L_{xx}=h''Y_x^2+h'Y_{xx}=h'Y_{xx}$ at $Y_x=0$ and $dL^*/dc=L_c=h'Y_c$, with $h'<0$. At fixed $c$, $dn=-dx/2$, so
$L_{nn}|_C=4L_{xx}$. For a monotone $g$, $(g\circ L)_{xx}=g''L_x^2+g'L_{xx}=g'L_{xx}$ at $L_x=0$ and
$d(g\circ L^*)/dc=g'\,dL^*/dc$; away from the argmin the term $g''L_x^2$ does not vanish.
(iii) $\partial_x\ln w=(Y_{cx}-Y_{xx})/(Y_c-Y_x)-(Y_{cx}+Y_{xx})/(Y_c+Y_x)$, which is $-2Y_{xx}/Y_c$ at $Y_x=0$.
(iv) Second-order chain rule. (v) Along the path $y^*=\gamma(c-\ln6)-\ln K$ (Lemma~\ref{lem:alloc} applied to the inner
aggregator), so $dy^*/dc=\gamma$; by (ii) and $1/\sigma^*-1=S/2$, $-Y_{xx}=\gamma S/4$ and $-y_{nn}|_C=-4Y_{xx}=S\gamma$. With
$\kappa=1$, $\gamma=abS$. (vi) Part (ii) uses only derivatives of a profile at its minimum and of the frontier; part (i)
uses two directional derivatives at the point.
\end{proof}

\begin{remark}[Finite grids]\label{rem:A-grid}
IsoFLOP profiles are sampled on finite grids and fitted by low-order polynomials \citep{czech2026problems}. If the profile
$f(t)$, with $t$ the distance from its minimum in $\ln N$ (or $\ln M$), is sampled at points symmetric about the minimum
and fitted by least squares on $(1,t,t^2)$, the quadratic coefficient is $f''/2+(f''''/24)\kappa_4+O(h^4)$, with
$\kappa_4\equiv(m_6-m_2m_4)/(m_4-m_2^2)$, $m_k$ the $k$-th moment of the grid and $h$ its half-width, and the fitted minimum
is displaced by $-(f'''/6)(m_4/m_2)/f''+O(h^3)$. For the \citet{besiroglu2024chinchilla} technology at $10^{21}$ FLOP and nine
equally spaced sizes, a grid spanning a factor of 9 (100) in $N$ overstates the curvature by 1.3 (5.9) percent, so the
model-free estimate is $\hat\sigma^*=0.734$ (0.726) against 0.737; a quartic fit removes the leading term (0.737 in both
cases). Quadratic fits on grids narrower than a factor of about 3 in $N$, or local polynomials of order four, keep the bias
below 0.001. These are noise-free calculations of bias only: with seed and evaluation noise, a narrower grid or a quartic
fit raises the variance of the estimated curvature, so the choice is a bias--variance trade-off to be made with the design's
noise level.
\end{remark}

\subsection{Partial identification beyond the design}\label{app:pi}

\begin{proposition}[Partial identification of $M^*(C)$ and of the wedge; Proposition~\ref{prop:pi} in the main text]\label{prop:A-pi}
Let $x^*(c)\equiv\ln M^*(C)$ be the unique argmin of the IsoFLOP profile at compute $C$, absolutely continuous in $c$, with
path slope $e(c)\equiv dx^*/dc=1-2a(c)$, where $a(c)\equiv d\ln N^*/d\ln C$. Suppose $x^*(c_J)$ is identified at the largest
design budget $c_J$ (from IsoFLOP minima, or from a path fitted to the minima at the largest budgets), and consider a model
with compute $c>c_J$ and log ratio $x=\ln M$.
\begin{enumerate}
\item[(i)] (Normal inputs.) If $N^*(C)$ and $D^*(C)$ are nondecreasing, then $e\in[-1,1]$ and
$x^*(c)\in[x^*(c_J)-(c-c_J),\ x^*(c_J)+(c-c_J)]$.
\item[(ii)] (Bounded path slope.) If $e(c')\in[e_L,e_U]$ for $c'\in[c_J,c]$, the sharp identified set is
\[
\mathcal X^*(c)=\big[x^*(c_J)+e_L(c-c_J),\ x^*(c_J)+e_U(c-c_J)\big],
\]
whose width grows linearly in the log extrapolation distance; monotonicity of $M^*$ is the case $e_L\ge0$. When the true
technology is in the generalized family, every point of $\mathcal X^*(c)$ is generated by a technology with positive
marginal products and single-peaked IsoFLOP profiles that coincides with the true one for $c\le c_J$. If confidence
intervals $\mathrm{CI}_j$ for $x^*(c_j)$ at budgets $c_1<\dots<c_J$ hold jointly with probability $1-p$ and the slope bounds
hold from $c_1$ on, then $\bigcap_j[\inf\mathrm{CI}_j+e_L(c-c_j),\ \sup\mathrm{CI}_j+e_U(c-c_j)]$ covers $x^*(c)$ with
probability at least $1-p$.
\item[(iii)] (Sign of the wedge.) If each IsoFLOP profile is single-peaked in $x$, which holds whenever $0<\sigma<1$
throughout (Lemma~\ref{lem:soc}), then $\operatorname{sign}(w-1)=\operatorname{sign}\big(x-x^*(c)\big)$ at every $(c,x)$, for
any technology. Hence $w>1$ is identified if and only if $x>\sup\mathcal X^*(c)$, $w<1$ if and only if
$x<\inf\mathcal X^*(c)$, and the sign is not identified otherwise. Neither $\sigma^*$ nor the functional form is needed.
\item[(iv)] (Magnitude.) Suppose $\partial\ln w/\partial x|_c\in[k_L,k_U]$ with $k_L>0$ between $x^*(c)$ and $x$; in the
generalized family $k_L=k_U=1/\sigma^*-1$, so a range for $\sigma^*$ gives $[k_L,k_U]$. With
$\mathcal X^*(c)=[\ell,u]$: $\ln w\in[k_L(x-u),\,k_U(x-\ell)]$ if $x>u$; $\ln w\in[k_U(x-u),\,k_L(x-\ell)]$ if $x<\ell$; and
$\ln w\in[k_U(x-u),\,k_U(x-\ell)]$ if $\ell\le x\le u$.
\item[(v)] (Expenditure share and $T$.) Since $s=1-1/w$ is increasing in $w$, bounds $[w_L,w_U]$ give
$s\in[1-1/w_L,\,1-1/w_U]$. Under Proposition~\ref{prop:A-wedge}(iii)(a) with $\phi=1$ and $\eta=\delta=0$,
$\lambda pT/(3D)\in[w_L-1,\,w_U-1]$; under the other objectives of Proposition~\ref{prop:A-wedge} the same bounds apply to
$w$ and map into $T$ as described there.
\item[(vi)] (A bound inside the support.) Let $F_{nn},F_{dd}\le0$ and $F_{nd}\ge0$, which hold throughout the generalized
family. Then $\ln w$ is nonincreasing in $n$ and nondecreasing in $d$, so for any support point with $n_b\ge n_0$ and
$d_b\le d_0$, $w(n_0,d_0)\ge w(n_b,d_b)$, and under (iii)(a) with $\lambda=p=1$,
$T_0\ge3D_0[w(n_b,d_b)-1]$, where $w(n_b,d_b)$ is identified without extrapolation. No finite upper bound follows from
these shape restrictions, and without $F_{nd}\ge0$ the bound can fail.
\end{enumerate}
\end{proposition}

\begin{proof}
(i) $d\ln N^*+d\ln D^*=d\ln C$ with both terms nonnegative gives $a\in[0,1]$; integrate $e=1-2a$. (ii) Integrate $e$ over
$[c_J,c]$; the linear continuations $x^*(c_J)+e(c-c_J)$, $e\in[e_L,e_U]$, attain every point. For the technology: in the
generalized family let $\tilde L(n,d)=E+[Ae^{-a_1(n+\psi_N(c))}+Be^{-b_1(d+\psi_D(c))}]^\kappa$ with $c=\ln6+n+d$ and
nondecreasing $\psi_N,\psi_D$ that vanish for $c\le c_J$. At fixed $c$ the profile is a translate of a profile of the
original technology (at effective compute $c+\psi_N+\psi_D$), hence single-peaked, with argmin $x^*(c)+2b\psi_N(c)-2a\psi_D(c)$ (Lemma~\ref{lem:alloc}(iii));
$\tilde F_n=F_n(1+\psi_N')+F_d\psi_D'>0$ and $\tilde F_d=F_n\psi_N'+F_d(1+\psi_D')>0$. Choosing $\psi_N'=\Delta'_+/(2b)$ and
$\psi_D'=\Delta'_-/(2a)$ generates any continuation $x^*(c)+\Delta(c)$ with $\Delta(c_J)=0$. On the event that every
$\mathrm{CI}_j$ covers, every cone contains $x^*(c)$. (iii) By Proposition~\ref{prop:A-modelfree}(i),
$w-1=-2Y_x/(Y_c+Y_x)=-2Y_x/F_d$ with $F_d>0$, and single-peakedness means $Y_x>0$ for $x<x^*$ and $Y_x<0$ for $x>x^*$. If
$0<\sigma<1$ throughout, every critical point of a profile is a strict local maximum of $y$
($Y_{xx}=-Q/(4F_n^2)<0$ at $Y_x=0$, Lemma~\ref{lem:soc}), so there is at most one. (iv) $\ln w(x)=\int_{x^*}^x\partial_x\ln w$,
with $\ln w(x^*)=0$ and $x^*\in[\ell,u]$. (v) Monotonicity. (vi) $\partial\ln w/\partial n=F_{nn}/F_n-F_{nd}/F_d\le0$ and
$\partial\ln w/\partial d=F_{nd}/F_n-F_{dd}/F_d\ge0$; move from $(n_b,d_b)$ to $(n_b,d_0)$ and then to $(n_0,d_0)$. In the
generalized family $F=-\kappa\ln Q$, whose Hessian has entries $a_1^2uv/Q^2$, $-a_1b_1uv/Q^2$ and $b_1^2uv/Q^2$ times
$-\kappa$. A saturating continuation of the data term ($\varepsilon_D\to0$ beyond the support) satisfies the shape
restrictions and sends $w\to\infty$; the concave quadratic $F=0.05n+0.3d-0.025n^2-0.005d^2-0.02nd$ has $F_{nd}<0$ and
$w(0,0.5)<w(0,0)$.
\end{proof}

\begin{remark}[An illustration at frontier scale]\label{rem:A-pi-illus}
This remark illustrates the logic of Proposition~\ref{prop:A-pi} only; the application in Section~\ref{sec:wedge}
harmonizes units and uses lab-own anchors. Anchor the path at the largest budgets of the two lab designs with the most
compute: the Chinchilla IsoFLOP minima (nine budgets up to $2.9\times10^{21}$ FLOP) and the Llama~3 minima (ten budgets up
to $10^{22}$ FLOP). A line fitted to each set of minima gives $M^*=20.5$ (95 percent interval $[11.4,36.8]$) and $22.3$
($[12.2,40.2]$) at the largest budget, with in-sample path slopes $e=0.005$ ($[-0.16,0.17]$) and $0.074$ ($[-0.09,0.24]$).
The intervals are equal-tailed wild bootstrap-$t$ intervals conditional on the designed budgets (HC2-rescaled residuals and
HC2 standard errors, Webb weights, 9,999 draws); with nine or ten points and the anchor at the high-leverage end of the
design, even this construction covers only about 90 percent in simulations at these designs, and a straight path within
the design is maintained. With the path slope confined to the range of allocation exponents across the paper's sweeps,
$a\in[0.37,0.57]$ or $e\in[-0.14,0.26]$, the identified set for $M^*(10^{24})$ is $[5.0,168]$ (Chinchilla anchor) or
$[6.4,133]$ (Llama~3 anchor). Among the verified-sample models beyond each anchor's largest budget (149 and 137 models),
the sign of $w-1$ is identified as positive for 115 and 111, negative for 2 and 2, and not identified for 32 and 24. With
only normal inputs, positive signs are identified for 46 and 47 models, and none is identified as negative. Point
extrapolation classifies every model, and the two anchors then disagree about the largest one (Llama~3.1 405B: $w>1$ with
the Chinchilla anchor, $w<1$ with the Llama~3 anchor), whose sign no set-valued assumption identifies. Magnitudes are
wide: with $\sigma^*\in[0.6,0.75]$ and the cross-sweep slope range, the medians of the models' lower and upper bounds on
$w$ are 1.58 and 16.1 (Chinchilla anchor). The models are measured in their own parameter and token conventions, which are
not harmonized with the anchors'.
\end{remark}

\begin{remark}[Parametric projection, calibration and the direction of extrapolation errors]\label{rem:A-pi-param}
A confidence set for $(\ln(A/B),\alpha,\beta)$ projects into one for $T$. At fixed compute,
$\Var(\ln\hat w)=\tfrac14\Var(\hat S)(\ln M)^2-\tfrac12\Cov(\hat S,\hat S\ln\hat M^*)\ln M+\tfrac14\Var(\hat S\ln\hat M^*)$ is
exactly quadratic in $\ln M$, so the band widens linearly in the log extrapolation distance with slope
$\tfrac12z_{1-p/2}\,\mathrm{sd}(\hat S)$. For the Chinchilla extraction, a pairs bootstrap ($B=300$) gives
$\mathrm{sd}(\ln\hat w)=0.161$ at $M=1{,}875$ and $C=7.2\times10^{23}$, of which 0.143 comes from $\ln M^*$ and 0.047 from
$\alpha+\beta$. That is sampling error within one sweep and one functional form ($\kappa=1$): the implied 95 percent
interval for $\ln M^*$ at that compute is about 1.6 log points wide, while the identified set of
Remark~\ref{rem:A-pi-illus} at $10^{24}$ FLOP is about 3.5 log points wide and the estimates of $M^*(10^{24})$ across the
fifteen technology estimates used for the wedge range from 1.4 to 291 (5.4 log points); the range of $\sigma^*$ adds more. Normalizing a family's flagship to $w_f=1$ is a calibration assumption, not the approach of
\citet{gandhi2020identification}, which uses an observed expenditure share; by Proposition~\ref{prop:A-wedge} it is valid
only if the flagship's $m_N$, $m_D$, $\delta$ and binding constraints all vanish. Finally, evidence that fitted laws
misstate losses at extreme ratios \citep{sardana2024beyond} does not sign the error in $\hat w$: revealed preference recovers
the technology the developer believed when it chose $(N,D)$ (Proposition~\ref{prop:A-wedge}(ii)), so if developers planned
with fitted laws, correcting toward the true technology moves $\hat w$ away from their first-order condition.
\end{remark}

\subsection{Growth-accounting corollaries}\label{app:growth}

\begin{corollary}[Frontier-release regressions]\label{cor:sahal}
Let frontier releases satisfy $\ln R_t=\ln K-\gamma(c_t+g_At)$, where $g_A$ is algorithmic progress in compute-equivalent
units, and $c_t=c_0+gt+u_t$ with $u_t$ independent of $t$. The OLS slope of $-\ln R_t$ on $c_t$ converges to
$\gamma[1+(g_A/g)R^2_{c,t}]$, with $R^2_{c,t}=g^2\Var(t)/(g^2\Var(t)+\Var(u))$. With deterministic trends this is
$\gamma/(1-s_A)$, $s_A\equiv g_A/(g+g_A)$: for $s_A\in[0.05,0.40]$ the naive scaling exponent is inflated by a factor of
1.05--1.67 \citep{nagy2013statistical}. Adding a time trend recovers $\gamma$ from deviations of compute from trend, and
does so only if $\Var(u)>0$.
\end{corollary}

\begin{proof}
$\Cov(-\ln R_t,c_t)=\gamma\Var(c)+\gamma g_A\Cov(t,c)$ and $\Cov(t,c)=g\Var(t)$. The last statement is Frisch--Waugh.
\end{proof}

\begin{corollary}[Hall-type bias in cost-share growth accounting]\label{cor:hall}
To first order, $\Delta y=\varepsilon_N\Delta n+\varepsilon_D\Delta d+\Delta\Omega$. A residual computed with equal weights
$\bar\varepsilon=(\varepsilon_N+\varepsilon_D)/2$, as the compute accounting $C=6ND$ suggests, equals
\[
\Delta\Omega+\tfrac12(\varepsilon_N-\varepsilon_D)(\Delta n-\Delta d)=\Delta\Omega+\bar\varepsilon\,\frac{w-1}{w+1}\,(\Delta n-\Delta d).
\]
Over-trained models ($w>1$) that shift further toward tokens ($\Delta d>\Delta n$) make the residual understate progress
\citep[the analog of][]{hall1988relation}. With weights $\gamma$ on each input, the bias is
$(\varepsilon_N-\gamma)\Delta n+(\varepsilon_D-\gamma)\Delta d$.
\end{corollary}

\begin{proof}
Subtract $\bar\varepsilon(\Delta n+\Delta d)$ and use $\varepsilon_N=2\bar\varepsilon w/(1+w)$ and
$\varepsilon_D=2\bar\varepsilon/(1+w)$. For discrete changes the error is second order; Törnqvist (average-of-endpoint)
weights reduce it to third order.
\end{proof}
